\documentclass[12pt]{article}
\usepackage[margin=1in]{geometry}
\usepackage{amsmath,amssymb,amsthm,mathtools}
\usepackage{physics}
\usepackage{booktabs}
\usepackage{array}
\usepackage{enumitem}
\usepackage[colorlinks=true,linkcolor=blue,citecolor=blue,urlcolor=blue]{hyperref}
\usepackage{cleveref}
\usepackage{cite}
\usepackage{graphicx}
\usepackage{xcolor}
\usepackage{microtype}

%--- Theorem environments ---
\newtheorem{theorem}{Theorem}[section]
\newtheorem{proposition}[theorem]{Proposition}
\newtheorem{definition}[theorem]{Definition}
\newtheorem{corollary}[theorem]{Corollary}
\newtheorem{remark}[theorem]{Remark}

%--- Shorthand macros ---
\newcommand{\FF}{\mathbb{F}}
\newcommand{\CC}{\mathbb{C}}
\newcommand{\ZZ}{\mathbb{Z}}
\newcommand{\Sp}{\mathrm{Sp}}
\newcommand{\PSp}{\mathrm{PSp}}
\newcommand{\SRG}[4]{\mathrm{SRG}(#1,#2,#3,#4)}
\newcommand{\GF}[1]{\mathbb{F}_{#1}}
\newcommand{\PG}[2]{\mathrm{PG}(#1,#2)}
\newcommand{\AG}[2]{\mathrm{AG}(#1,#2)}
\newcommand{\Aut}{\mathrm{Aut}}

\title{\textbf{Universal Quantum Computation Through a Single Photon}\\[6pt]
\large Symplectic Polar Geometry, the Ternary Field $\GF{3}$, Qutrits,\\
and the Photonic Gate Universe of $W(3,3)$}

\author{%
  Wil Dahn\\[4pt]
  \normalsize\textit{Co-authored in partnership with the DTR \& UOR Foundation}\\[2pt]
  \normalsize\textit{Theory of Everything Project}%
}

\date{May 2026}

\begin{document}
\maketitle

\begin{abstract}
We present a unified account of how a single photon---the most elementary quantum of
the electromagnetic field---can serve as the physical carrier for universal quantum
computation, grounded in the rich combinatorial and algebraic structure of the
symplectic polar space $W(3,3)$.  The space $W(3,3)$ is the incidence geometry of
totally isotropic 1-spaces in $\PG{3}{\GF{3}}$ with respect to a non-degenerate
alternating bilinear form; its collinearity graph is the unique strongly regular
graph $\SRG{40}{12}{2}{4}$.  Crucially, $W(3,3)$ is built on the \emph{ternary
field} $\GF{3} = \{0,1,2\}$, which encodes a three-level quantum system---a
\emph{qutrit}.  The field $\GF{3}$ is \emph{uniquely} selected by the Diophantine
equation $q! = 2q$, whose only non-trivial solution is $q = 3$; from this single
integer the entire parameter table of $W(3,3)$---vertices $v=40$, valency $k=12$,
intersection parameters $\lambda=2$, $\mu=4$, and automorphism group
$\Sp(4,\FF_3)$ of order $51840$---follows by closed arithmetic.  We show that the
polarisation degrees of freedom of a single photon implement the qubit face of this
qutrit geometry ($\lambda = 2$), that three orthogonal optical modes realise a full
photonic qutrit, and that the twelve-element valency $k=12$ captures the generator
count of the two-qutrit Clifford group $\Sp(4,\FF_3)$.  Combined with the
Knill--Laflamme--Milburn (KLM) protocol and measurement-based quantum computing
(MBQC), we obtain a complete framework in which all cardinal numbers of photonic
universality are read directly from the $W(3,3)$ parameter table.
\end{abstract}

%----------------------------------------------------------------------
\section{Introduction}
%----------------------------------------------------------------------

Quantum computing promises exponential speedups for a select but important class of
problems, yet its physical realisation depends crucially on identifying quantum systems
that are both controllable and scalable.  Among all candidate platforms, \emph{linear
optical quantum computing} (LOQC) occupies a distinguished position: it operates at
room temperature, suffers no decoherence from phonons or nuclear spins, and
communicates at the speed of light~\cite{KLM2001}.

The fundamental unit of LOQC is the \emph{single photon}.  A photon is a Fock-state
excitation of the quantised electromagnetic field,
\begin{equation}
  \ket{1} = \hat{a}^\dagger \ket{0},
  \label{eq:fock}
\end{equation}
where $\hat{a}^\dagger$ is the creation operator satisfying
$[\hat{a},\hat{a}^\dagger]=1$.  Despite being a single quantum, the photon possesses
multiple independent degrees of freedom---polarisation, path (spatial mode), time-bin,
and orbital angular momentum (OAM)---each of which can encode information.  The
natural encoding is a \emph{qubit} ($\CC^2$) for two-dimensional subspaces; but OAM
ladder modes and three-mode architectures elevate the photon to a \emph{qutrit}
($\CC^3$), unlocking a richer computational and error-correcting landscape.

What determines which dimensions are special?  We argue that the answer is the
symplectic polar space $W(3,3)$, the unique geometry derived from the ternary field
$\GF{3}$ and a non-degenerate alternating form on $\GF{3}^4$.  As we explain in
\cref{sec:w33}, this geometry is not chosen by hand: it is \emph{forced} by the
integer equation $q! = 2q$, which has the unique non-trivial solution $q = 3$.
All structural constants of single-photon computation---the qubit dimension~$2$, the
qutrit dimension~$3$, the spacetime dimension~$4$, and the twelve generators of
$\Sp(4,\FF_3)$---read off directly from the $W(3,3)$ parameter table.

This paper is structured as follows.  \Cref{sec:w33} introduces $W(3,3)$ from first
principles: the symplectic polar space definition, the Diophantine seed, and the
closed-form parameter derivation.  \Cref{sec:qutrits} establishes the connection
between $\GF{3}$ and three-level quantum systems.  \Cref{sec:hilbert} reviews the
physical Hilbert space of a single photon and its qutrit extension.
\Cref{sec:gates} presents universal gate sets for qubits and qutrits.
\Cref{sec:klm} describes the KLM protocol for two-qubit gates.  \Cref{sec:mbqc}
connects single-photon sources to MBQC via cluster states.  \Cref{sec:runtime}
separates the probabilistic, deterministic, quantum, and classical runtime layers.
\Cref{sec:encoding} presents the structural encoding theorem, and
\cref{sec:conclusions} concludes.

%======================================================================
\section{The Symplectic Polar Space \texorpdfstring{$W(3,3)$}{W(3,3)}}
\label{sec:w33}
%======================================================================

Before we can say anything meaningful about the role of $W(3,3)$ in quantum
computation, we must define it properly.  $W(3,3)$ is \emph{not} merely a graph with
four parameters: it is a rich incidence geometry, a classical object in finite
geometry, whose automorphism group connects to the exceptional Lie algebra $E_6$ and
the two-qutrit Clifford group simultaneously.

\subsection{The Master Equation: Why \texorpdfstring{$q=3$}{q=3}}

\begin{theorem}[Uniqueness of the ternary field]\label{thm:dioph}
The equation
\begin{equation}
  q! \;=\; 2q
  \label{eq:dioph}
\end{equation}
has a unique non-trivial solution in the positive integers: $q = 3$.
\end{theorem}

\begin{proof}
For $q=1$: $1!=1 \neq 2$.  For $q=2$: $2!=2 \neq 4$.  For $q=3$:
$3!=6=2\cdot 3$.\ \checkmark.  For $q\ge 4$: $q! \ge 24 > 2q$, and the factorial
grows super-exponentially while $2q$ grows linearly; no further solution exists.
\end{proof}

The number $q=3$ is the \emph{only} prime power for which the factorial coincides
with twice itself.  We take this as the canonical selector for the underlying field:
the geometry is built over $\GF{3}$, the field with three elements.

\subsection{The Symplectic Polar Space}

\begin{definition}[Symplectic polar space $W(2n-1,q)$]\label{def:sympl}
Let $V = \GF{q}^{2n}$ and let $\beta : V \times V \to \GF{q}$ be a non-degenerate
alternating bilinear form (i.e., $\beta(u,u) = 0$ for all $u \in V$ and
$\beta(u,v) = -\beta(v,u)$).  A subspace $W \le V$ is \emph{totally isotropic} if
$\beta(u,v) = 0$ for all $u,v \in W$.  The \textbf{symplectic polar space}
$W(2n-1,q)$ is the incidence geometry whose:
\begin{itemize}[noitemsep]
  \item \emph{points} are the totally isotropic 1-spaces (projective points) of
    $\PG{2n-1}{\GF{q}}$;
  \item \emph{lines} are the totally isotropic 2-spaces (projective lines) of
    $\PG{2n-1}{\GF{q}}$.
\end{itemize}
\end{definition}

For $n=2$ and $q=3$ we obtain $W(3,3)$: the geometry of totally isotropic
1-spaces of $\GF{3}^4$.  The canonical alternating form is
\begin{equation}
  \beta\bigl((u_1,u_2,u_3,u_4),\,(v_1,v_2,v_3,v_4)\bigr)
  \;=\; u_1 v_3 - u_3 v_1 + u_2 v_4 - u_4 v_2,
\end{equation}
with coefficients in $\GF{3}$.

\subsection{Points, Lines, and the Collinearity Graph}

\begin{proposition}[Points and lines of $W(3,3)$]\label{prop:pts}
The symplectic polar space $W(3,3)$ has:
\begin{itemize}[noitemsep]
  \item \textbf{40 points}: the totally isotropic 1-spaces of $\GF{3}^4$.  The count
    is
    \begin{equation}
      \frac{3^4 - 1}{3 - 1} \;=\; \frac{80}{2} \;=\; 40.
      \label{eq:pts}
    \end{equation}
  \item \textbf{40 lines}: totally isotropic 2-spaces, each containing $q+1=4$ points.
  \item \textbf{Collinearity graph}: the strongly regular graph $\SRG{40}{12}{2}{4}$,
    unique up to isomorphism.
\end{itemize}
\end{proposition}

\subsection{Parameter Derivation from $q=3$}

\begin{definition}[SRG eigenvalue quadratic]\label{def:srg}
A strongly regular graph $\SRG{v}{k}{\lambda}{\mu}$ has the property that every pair
of adjacent vertices shares exactly $\lambda$ common neighbours, and every pair of
non-adjacent vertices shares exactly $\mu$.  The restricted eigenvalues $r, s$ satisfy
\begin{equation}
  x^2 - (\lambda - \mu)\,x - (k - \mu) \;=\; 0.
  \label{eq:srg_quad}
\end{equation}
\end{definition}

\begin{proposition}[Parameter closure from $q=3$]\label{prop:params}
Setting $\lambda = q-1 = 2$ and $\mu = q+1 = 4$ in \cref{eq:srg_quad} yields
$x^2 + 2x - (k-4) = 0$, with roots $x = -4$ and $x = 2$ when $k = 12$.  The
feasibility identity $k(k-\lambda-1) = (v-k-1)\mu$ then forces $v = 40$.  The full
parameter table is:
\[
  v = 40,\quad k = 12,\quad \lambda = 2,\quad \mu = 4,
\]
\[
  r = 2 \;(f = 24), \quad s = -4 \;(g = 15), \quad E = 240.
\]
The automorphism group satisfies $|\Aut(W(3,3))| = |\Sp(4,\FF_3)| = 51840$.
\end{proposition}

\begin{remark}[Trit-savings principle]
\Cref{eq:pts} encodes a \emph{trit-savings} calculation: $3^4 = 81$ exponent vectors
in $\GF{3}^4$, minus the zero vector, divided by the two non-zero scalars
$\GF{3}^\times = \{1,2\}$, collapses to exactly $80/2 = 40$ projective observables.
This is the three-state accounting that makes $W(3,3)$ the natural two-qutrit phase
space.
\end{remark}

\subsection{Exceptional Isomorphism and the Weyl Group of \texorpdfstring{$E_6$}{E6}}

The automorphism group of $W(3,3)$ enjoys the exceptional isomorphism
\begin{equation}
  \PSp(4,3) \;\cong\; W(E_6)^+,
\end{equation}
where $W(E_6)^+$ is the index-2 positive Weyl group of the exceptional Lie algebra
$E_6$, of order $25920$.  The full automorphism group $\Sp(4,\FF_3)$ of order
$51840 = 2 \cdot 25920$ is the full Weyl group $W(E_6)$.  The same $q=3$ seed also
generates the exceptional Lie series arithmetically:
\begin{equation}
  G_2 = k + \lambda = 14, \quad
  F_4 = \mu\,\Phi_3 = 52, \quad
  E_6 = \lambda\,q\,\Phi_3 = 78, \quad
  E_7 = \Phi_3\Phi_4 + q = 133, \quad
  E_8 = E + \lambda^q = 248,
\end{equation}
where $\Phi_3 = 13$ and $\Phi_4 = 10$ are cyclotomic-polynomial values at $q=3$.

%======================================================================
\section{The Ternary Field \texorpdfstring{$\GF{3}$}{F3} and Qutrits}
\label{sec:qutrits}
%======================================================================

\subsection{From a Field to a Hilbert Space}

The field $\GF{3} = \{0, 1, 2\}$ (equivalently $\ZZ/3\ZZ = \{0, 1, -1\}$) has
exactly three elements.  In quantum information, a \emph{qutrit} is a quantum system
whose Hilbert space is $\CC^3$, spanned by three orthonormal basis states
$\{\ket{0},\ket{1},\ket{2}\}$.  A pure qutrit state is
\begin{equation}
  \ket{\psi} = \alpha_0\ket{0} + \alpha_1\ket{1} + \alpha_2\ket{2}, \qquad
  \sum_{j=0}^{2}|\alpha_j|^2 = 1.
\end{equation}
The correspondence is immediate: the three elements of $\GF{3}$ label the three
computational basis states of a qutrit.  The ternary field is the \emph{index set}
of the qutrit Hilbert space, and \cref{thm:dioph} shows this choice is unique.

\begin{definition}[Two-qutrit Pauli group]\label{def:pauli3}
The generalised Pauli operators on $\CC^3$ are generated by the clock operator $Z_3$
and shift operator $X_3$:
\begin{equation}
  Z_3\ket{j} = \omega^j\ket{j}, \qquad X_3\ket{j} = \ket{j+1 \bmod 3},
  \qquad \omega = e^{2\pi i/3}.
\end{equation}
For two qutrits, the $3^4 = 81$ Pauli monomials
$X_3^{a_1}Z_3^{b_1} \otimes X_3^{a_2}Z_3^{b_2}$ (with $(a_1,b_1,a_2,b_2)\in\GF{3}^4$)
generate the two-qutrit Pauli group.  Modulo global phases and the zero vector, this
gives exactly the \textbf{40 projective observables} of \cref{prop:pts}.
\end{definition}

\begin{theorem}[$W(3,3)$ as two-qutrit phase space]\label{thm:phasespace}
The 40 points of $W(3,3)$ are in canonical bijection with the non-zero two-qutrit
Pauli observables (up to phase).  Adjacency in $W(3,3)$ corresponds to mutual
\emph{commutativity} of the Pauli operators.  The automorphism group $\Sp(4,\FF_3)$
is precisely the two-qutrit Clifford group modulo phases.
\end{theorem}

\begin{proof}[Proof sketch]
A non-zero vector $(a_1,b_1,a_2,b_2) \in \GF{3}^4$ determines a Pauli monomial; two
monomials commute iff their symplectic inner product
$a_1 b_1' - b_1 a_1' + a_2 b_2' - b_2 a_2'$ vanishes mod $3$, i.e., the
corresponding projective vectors are adjacent in $W(3,3)$.  The group of symplectic
linear maps $\Sp(4,\FF_3)$ acts faithfully on this phase space.
\end{proof}

\subsection{Why Three Levels Are Computationally Special}

\begin{enumerate}[label=(\alph*)]
  \item \textbf{Magic state universality.}  For qutrits, the cubic phase gate
    $T_3 = \mathrm{diag}(1, \omega, \omega^4)$ provides magic injection.  With
    Clifford$+T_3$, the two-qutrit gate set is universal (Bravyi--Kitaev~\cite{BK2005}).
    The magic state derives from the cubic invariant of $\GF{3}^{1+2}$.
  \item \textbf{Trit-savings.}  $3^4 = 81$ exponent vectors collapse to
    $80/2 = 40$ projective observables (cf.\ \cref{eq:pts}).
  \item \textbf{Minimal universal Turing machine.}  The smallest confirmed universal
    Turing machine uses $(\lambda, q) = (2, 3)$ states and symbols
    (Wolfram--Smith~\cite{Smith2007}).  Both $\lambda = 2$ and $q = 3$ are $W(3,3)$
    parameters.
\end{enumerate}

\subsection{The Two-Qutrit Clifford Group as a Gate Set}

The group $\Sp(4,\FF_3)$ of order $51840$ is the exact automorphism group of the
two-qutrit commutation relations.  Its generators include:
\begin{itemize}[noitemsep]
  \item Hadamard-like Fourier gates $F_3$ on each qutrit;
  \item Controlled-$Z$ gate $CZ_3$ between the two qutrits;
  \item Phase gates $S_3$.
\end{itemize}
The valency $k = 12$ of $W(3,3)$ counts the generators of the two-qutrit Clifford
algebra.  Augmented by $T_3$, the gate set $\{F_3, CZ_3, S_3, T_3\}$ is
\emph{universal} for quantum computation on qutrit registers.

%======================================================================
\section{The Single-Photon Hilbert Space}
\label{sec:hilbert}
%======================================================================

\subsection{Polarisation Qubit: The \texorpdfstring{$\lambda=2$}{lambda=2} Face of \texorpdfstring{$W(3,3)$}{W(3,3)}}

The polarisation degree of freedom of a single photon spans a two-dimensional Hilbert
space $\mathcal{H}_{\mathrm{pol}} \cong \mathbb{C}^2$:
\begin{equation}
  \ket{\psi} = \alpha\ket{H} + \beta\ket{V}, \quad |\alpha|^2+|\beta|^2 = 1,
\end{equation}
where $\ket{H}$ and $\ket{V}$ denote horizontal and vertical polarisations.  The
dimension~$2$ is precisely the parameter $\lambda = 2$ of $W(3,3)$: the two-fold
degeneracy of photon spin mirrors the two-fold common-neighbour count of adjacent
graph vertices.

\subsection{Dual-Rail (Path) Encoding}

Two spatial modes provide an alternative qubit encoding:
\begin{equation}
  \ket{0}_L = \hat{a}_0^\dagger\ket{0,0} = \ket{1,0}, \qquad
  \ket{1}_L = \hat{a}_1^\dagger\ket{0,0} = \ket{0,1}.
\end{equation}
Photon loss is heralded (vacuum in both modes), making this encoding inherently
error-detectable.

\subsection{Photonic Qutrit: Three Optical Modes}

The ternary field $\GF{3}$ underlying $W(3,3)$ motivates extending the photon's
computational space to a \emph{qutrit}.  Three physical realisations exist:

\begin{enumerate}[label=(\alph*)]
  \item \textbf{Orbital angular momentum (OAM).}  The three modes
    $\ell \in \{-1, 0, +1\}$ define an OAM qutrit.
    Laguerre--Gaussian beam optics provides the gate set; SU(3) rotations are
    implemented by spatial light modulators~\cite{Erhard2018}.
  \item \textbf{Three-path encoding.}  A single photon in three spatial modes
    $\hat{a}_j^\dagger\ket{0,0,0}$, $j\in\{0,1,2\}$, constitutes a path-encoded
    qutrit.
  \item \textbf{Polarisation$+$path hybrid.}  Combined modes
    $\{H_0, V_0, H_1\}$ implement the $\GF{3}$ label set $\{0,1,2\}$ directly.
\end{enumerate}

The three-mode photon is the \emph{physical realisation of a $W(3,3)$ point}: each
of the 40 points is a projective qutrit state in $\PG{3}{\GF{3}}$, and collinearity
is the physical commutation relation of the corresponding Pauli observables.

%======================================================================
\section{Self-Entanglement: The Temporal Qutrit on a Single Photon}
\label{sec:self-entanglement}
%======================================================================

The qutrit structure introduced in \cref{sec:qutrits,sec:hilbert} treats the
photon's three-level mode space as a \emph{spatial} or \emph{spectral} register.
We now show that a single photon also carries a \emph{temporal} qutrit---the three
modes $\{\text{past},\text{now},\text{future}\}$---and that the natural state on
this temporal register is the maximally self-entangled Bell-qutrit
\begin{equation}
  \ket{\Omega} \;=\; \frac{1}{\sqrt 3}\,\bigl(\ket{0}_p\ket{0}_f
    \;+\;\ket{1}_p\ket{1}_f \;+\;\ket{2}_p\ket{2}_f\bigr),
  \label{eq:bellqutrit}
\end{equation}
where the subscripts $p,f$ label past and future temporal modes.  The ``now'' mode
arises as the harmonic convergence of past and future under the Choi--Jamio\l kowski
identity; the resulting carrier is a single photon whose internal temporal qutrit
\emph{is} the substrate of $W(3,3)$ universal computation.

\subsection{Past/Future Qutrit Register on One Photon}

A single photon traversing a delay-line interferometer can occupy three
time-bin modes $\{t_0, t_1, t_2\}$ separated by a fixed delay $\tau$.  Label the
\emph{past} register $p$ by the early time-bin pair $\{t_0,t_1\}$ and the
\emph{future} register $f$ by the late pair $\{t_1,t_2\}$; the overlap at $t_1$
is the shared ``now'' bin.  Equivalently, one can use the two halves of an
asymmetric Mach--Zehnder interferometer, with the central recombination
implementing the now-register entanglement~\cite{Brendel1999,Marcikic2002}.

\begin{definition}[Temporal qutrit register]\label{def:tempqutrit}
The \emph{temporal qutrit} on a single photon is the three-dimensional Hilbert
space spanned by the time-bin basis $\{\ket{0},\ket{1},\ket{2}\}$ encoding
$\{\text{past},\text{now},\text{future}\}$.  A bipartite split into past and future
registers gives the product space
$\mathcal H_p \otimes \mathcal H_f = \CC^3 \otimes \CC^3 \cong \CC^9$,
with $q^2 = 9$ histories factoring as $q^2 = q + q!= 3 + 6$ (diagonal histories
plus off-diagonal context-changing histories).
\end{definition}

\subsection{The Bell Qutrit and the Choi--Jamio\l kowski Identity}

The maximally entangled state \eqref{eq:bellqutrit} satisfies the
Choi--Jamio\l kowski (CJ) identity: for every $U \in \mathrm{U}(3)$,
\begin{equation}
  \langle\Omega| (I_p\otimes U_f) |\Omega\rangle
  \;=\; \frac{\mathrm{Tr}(U)}{3}.
  \label{eq:choi}
\end{equation}
That is, applying \emph{any} unitary to the future register and projecting onto the
Bell state yields the unitary's normalised trace.  In particular, the photon's
temporal qutrit \emph{represents every unitary} by a single Bell-state overlap,
which is the operational meaning of the channel--state duality.

\begin{theorem}[Bell qutrit as Choi state of $W(3,3)$]\label{thm:bellchoi}
The Bell qutrit $\ket{\Omega}$ is the Choi state of the identity channel on a single
qutrit.  Its stabiliser group $\langle X_p X_f,\ Z_p Z_f^{-1}\rangle$ is a totally
isotropic 2-space of $\GF{3}^4$ under the canonical symplectic form, i.e.\ a
\emph{line of $W(3,3)$}.
\end{theorem}

\begin{proof}
The Choi state of a channel $\Phi$ is $(\Phi\otimes I)\ket{\Omega}\langle\Omega|$;
for $\Phi = I$, this is $\ket{\Omega}\langle\Omega|$ itself.  The operators
$X_p X_f$ and $Z_p Z_f^{-1}$ commute (their symplectic inner product is
$(1,0,1,0)\cdot J\cdot(0,1,0,-1)^T = 0$ over $\GF{3}$), and $\ket{\Omega}$ is the
unique simultaneous $+1$ eigenstate.  Two independent commuting Pauli generators
on $\GF{3}^4$ span a 2-dimensional isotropic subspace---by
\cref{prop:pts}, exactly a line of $W(3,3)$.
\end{proof}

\subsection{Now as the Harmonic Convergence}

The ``now'' mode is not an independent third register; it is the
\emph{harmonic convergence point} of the past/future Bell pair under the CJ
identity.  In time-bin optics, this is the recombination interferometer's central
output; in quantum-channel language, it is the channel-state correspondence point
where temporal evolution becomes spatial entanglement.

\begin{proposition}[Now as Choi convergence]\label{prop:now}
For any qutrit unitary $U$ acting on the future register, define the now-state
$\ket{\mathrm{now}_U} = (I_p\otimes U_f)\ket{\Omega}/\|\cdot\|$.  Then:
\begin{enumerate}[label=(\roman*),noitemsep]
  \item $\|(I_p\otimes U_f)\ket{\Omega}\|^2 = 1$ (norm-preserving),
  \item $\langle\mathrm{now}_U|\mathrm{now}_V\rangle = \mathrm{Tr}(U^\dagger V)/3$,
  \item the 9-dimensional space of now-states is exactly the Choi space of $\mathrm{U}(3)$.
\end{enumerate}
\end{proposition}

\subsection{Uniqueness: Now as the Forced Fixed Point}

The ``harmonic convergence'' interpretation of \cref{prop:now} is not a labelling
choice---the Bell qutrit $\ket{\Omega}$ is forced by symmetry.  The following
characterisation gives four equivalent conditions, each of which picks out
$\ket{\Omega}$ uniquely (up to a global phase) among all states of
$\mathcal H_p \otimes \mathcal H_f$.

\begin{theorem}[Uniqueness of the past/future Bell qutrit]\label{thm:bell_uniqueness}
Let $\mathcal H_p = \mathcal H_f = \CC^q$ with $q = 3$, and let
$\ket{\psi} \in \mathcal H_p \otimes \mathcal H_f$ be a normalised state.  The
following are equivalent (and each holds for the Bell qutrit $\ket{\Omega}$ of
\eqref{eq:bellqutrit} alone, up to a global phase):
\begin{enumerate}[label=(\roman*),noitemsep]
  \item \textbf{Time-reversal $+$ maximal entanglement:}
    $S_{pf}\ket{\psi} = \ket{\psi}$ (SWAP-invariance under past$\leftrightarrow$future
    exchange) AND the entanglement entropy is $S(\ket{\psi}) = \log q$.
  \item \textbf{Choi--Jamio\l kowski state of the identity channel:}
    $\ket{\psi} = (I_p\otimes I_f)\ket{\Omega}$.
  \item \textbf{Unitary--conjugate invariance:} $(U \otimes U^*)\ket{\psi} =
    \ket{\psi}$ for every $U \in \mathrm{U}(q)$.
  \item \textbf{Uniform Schmidt spectrum:} the Schmidt decomposition of
    $\ket{\psi}$ has all Schmidt coefficients equal to $1/\sqrt q$.
\end{enumerate}
\end{theorem}

\begin{proof}
(iv)$\Rightarrow$(ii):  A normalised state with all Schmidt coefficients equal
to $1/\sqrt q$ can be written as $\ket{\psi} = q^{-1/2}\sum_j \ket{\alpha_j}
\otimes\ket{\beta_j}$ for orthonormal Schmidt bases $\{\ket{\alpha_j}\},
\{\ket{\beta_j}\}$.  Changing basis to $\{\ket j\}$ on each register is a
unitary, so $\ket{\psi}$ and $\ket{\Omega}$ differ only by a unitary and a
global phase.

(ii)$\Rightarrow$(iii):  Direct calculation,
\begin{equation*}
  (U\otimes U^*)\ket{\Omega} \;=\; q^{-1/2}\sum_j U\ket{j}\otimes U^*\ket{j}
  \;=\; q^{-1/2}\sum_{a,b}\Bigl(\sum_j U_{aj}\overline{U}_{bj}\Bigr)\ket{ab}
  \;=\; q^{-1/2}\sum_{a,b}\delta_{ab}\ket{ab} \;=\; \ket{\Omega},
\end{equation*}
using $U^*\ket{j} = \sum_b\overline U_{bj}\ket{b}$ and unitarity
$\sum_j U_{aj}\overline U_{bj} = (UU^\dagger)_{ab} = \delta_{ab}$.

(iii)$\Rightarrow$(iv):  The decomposition of $\mathcal H_p\otimes\mathcal H_f$
under the action $(U\otimes U^*)$ is into irreducibles indexed by $\mathrm{SU}(q)$
representations; only the trivial representation contributes an invariant
vector, and that trivial subspace is one-dimensional (the symmetric maximally
mixed state's purification, equivalently $\ket{\Omega}$).  Hence any state fixed
by $(U\otimes U^*)$ for all $U$ is a scalar multiple of $\ket{\Omega}$, and its
Schmidt coefficients are uniformly $1/\sqrt q$.

(i)$\Leftrightarrow$(iv):  SWAP-invariance restricts $\ket{\psi}$ to a symmetric
tensor, $\ket{\psi} = \sum_{jk}c_{jk}\ket{jk}$ with $c_{jk}=c_{kj}$; entropy
$S(\ket{\psi})$ is maximised at $\log q$ exactly when the singular values of the
matrix $(c_{jk})$ are all $1/\sqrt q$, i.e.\ uniform Schmidt spectrum.

This closes (i)$\Leftrightarrow$(ii)$\Leftrightarrow$(iii)$\Leftrightarrow$(iv),
and the unique state (up to global phase) satisfying any of them is
$\ket{\Omega}$.
\end{proof}

\begin{corollary}[``Now'' is forced]\label{cor:now_forced}
The identification of the ``now'' state with the past/future Bell qutrit
$\ket{\Omega}$ is forced by either of two minimal physical requirements:
(a) time-reversal symmetry plus maximal past/future entanglement, or
(b) unitary-conjugate invariance $(U\otimes U^*)$.  ``Now $=$ harmonic
convergence'' is therefore not a phenomenological label but the unique fixed
point of these symmetries on the qutrit Hilbert space.
\end{corollary}

\begin{remark}[Physical realisation of the symmetries]\label{rem:phys_sym}
Time-reversal exchange corresponds to swapping early and late time-bins in the
Mach--Zehnder interferometer of \cref{def:tempqutrit}.  Unitary-conjugate
invariance corresponds to applying any $\mathrm{SU}(3)$ rotation jointly to the
past polarisation/path mode and the complex-conjugate ($\sigma_y$-flipped)
rotation to the future mode---the temporal analog of the spin-singlet rotation
symmetry of two qubits.  Both symmetries are \emph{a priori} required of any
consistent self-entangled temporal state on one photon (there is no preferred
temporal basis at the single-photon level), so \cref{thm:bell_uniqueness} forces
$\ket{\Omega}$ uniquely.
\end{remark}

\subsection{Spreads, Now-Contexts, and the 40-Ray Closure}

The Bell line of \cref{thm:bellchoi} is one of the 40 lines of $W(3,3)$.  Each line
participates in a \emph{spread} of $W(3,3)$ (a partition of the 40 points into
10 disjoint lines of 4 points each).  Each line in a spread is a
\emph{now-context}: a maximal set of mutually commuting Pauli observables that can
be jointly measured at a single instant.

\begin{theorem}[Now-context closure]\label{thm:nowclosure}
Every spread of $W(3,3)$ contains exactly $10$ disjoint now-contexts of $q+1=4$
rays each, giving the closure
\begin{equation}
  10 \cdot 4 \;=\; 40 \;=\; |W(3,3)|.
\end{equation}
The Bell line (\cref{thm:bellchoi}) lies in exactly $q^2 = 9$ spreads, and the
co-context cloud (the union of all lines incident to or disjoint from the Bell
line) closes as $1 + 12 + 27 = 40$ (Bell + intersecting + disjoint).
\end{theorem}

\begin{proof}[Proof sketch]
$W(3,3)$ admits $36 = (q!)^2 = \mu\cdot q^2$ spreads total; each line is in $q^2 = 9$ of them
(see~\cite{Vlasov2025}).  Within a spread, the lines are pairwise disjoint
maximal isotropic 2-spaces, so they form a partition of $40$ into $10$ blocks of
$4$ points each.  The Bell-local shell counts: $1$ (Bell line itself) +
$k = 12$ (lines meeting Bell at one point) + $27 = q^q$ (lines disjoint from Bell)
$= 40$.
\end{proof}

\subsection{Two Commuting Qutrits on One Photon: The Temporal--Spectral Pair}

A single photon naturally carries \emph{two} commuting qutrit registers when both a
temporal three-bin structure and a spectral three-sideband structure are
simultaneously addressed.  Following the photonic dictionary of
\cite[\S 3]{Vlasov2025}, the temporal register $\{p,n,f\} = \{\text{past},
\text{now}, \text{future}\}$ commutes with the spectral register
$\{\ell, c, u\} = \{\text{lower}, \text{carrier}, \text{upper}\}$ sideband.
Together they span a $3\times 3 = 9$-dimensional Hilbert space, isomorphic to
$\GF{3}^2$, which lifts to $\GF{3}^4$ (i.e.\ two stacked $\GF{3}^2$'s) by
including the qubit polarisation factor and the path-encoded mode---the full
$W(3,3)$ phase space.

\begin{proposition}[$3\times 3$ photonic torus]\label{prop:33torus}
The temporal $\otimes$ spectral two-qutrit register on a single photon is a
$3\times 3$ discrete torus.  Its $q^2+q+1 = 13$ projective points (the
\emph{screen}, equal to $\PG{2}{\GF{3}}$) form the \emph{memory layer}, and its
$q^q = 27$ affine points (the \emph{bulk}, equal to $\AG{3}{\GF{3}}$ minus the
plane at infinity) form the \emph{compute layer}.  These add to $13 + 27 = 40 =
|W(3,3)|$, recovering the full vertex count.
\end{proposition}

The $36 = (q!)^2$ complete two-qutrit measurement programs are exactly the
spreads of \cref{thm:nowclosure}, so every program in the temporal--spectral
photonic computer corresponds to a spread of $W(3,3)$.

\subsection{Self-Entanglement on a Single Photon}

We now state the main result of this section.

\begin{theorem}[Single-photon self-entanglement]\label{thm:selfent}
A single photon equipped with a temporal-qutrit register $\mathcal H_t = \CC^3$
admits a canonical self-entangled state $\ket{\Omega} \in \mathcal H_p \otimes
\mathcal H_f \subset \mathcal H_t^{\otimes 2}$ which:
\begin{enumerate}[label=(\roman*),noitemsep]
  \item is maximally entangled across the past/future split, with entanglement
    entropy $S(\ket{\Omega}) = \log q = \log 3$;
  \item has stabiliser equal to a line of $W(3,3)$ (\cref{thm:bellchoi});
  \item generates $40$ now-context rays through the spread of its stabiliser line
    (\cref{thm:nowclosure}, $10\cdot 4 = 40$);
  \item realises the Choi--Jamio\l kowski state of every $\mathrm{U}(3)$ channel
    via \eqref{eq:choi}.
\end{enumerate}
Consequently, one self-entangled qutrit on a single photon compiles to the entire
$W(3,3)$ universe via $(3^4 - 1)/(3-1) = 40$, Bell-cloud $1 + 12 + 27 = 40$, and
spread-frame closure $10 \cdot 4 = 40$.
\end{theorem}

\begin{proof}
(i) The Bell qutrit \eqref{eq:bellqutrit} has Schmidt decomposition with three
equal Schmidt coefficients $\lambda_j = 1/\sqrt 3$; entropy of entanglement
$S = -\sum \lambda_j^2 \log \lambda_j^2 = \log 3$.\\
(ii)--(iv) follow from \cref{thm:bellchoi,thm:nowclosure}, \eqref{eq:choi}.\\
The compilation identities $(3^4-1)/(3-1)=40$, $1+12+27=40$, $10\cdot 4=40$ all
match $|W(3,3)|$, so the single self-entangled qutrit compiles to the full
40-point $W(3,3)$ phase space.
\end{proof}

\begin{remark}[Minimality]
The self-entanglement is \emph{minimal}: for $q=2$ (qubit), the corresponding
substrate has only $v(2) = 15$ stabiliser points; for $q=3$ (qutrit), $v(3) = 40$.
The jump from $15$ to $40$ at $q=3$ is the Diophantine ``activation'' of universal
quantum computation on a single photon, consistent with the master equation
$q! = 2q$ (\cref{thm:dioph}).
\end{remark}

\subsection{Diagonal vs Off-Diagonal: The Master Equation on $\mathcal H_p\otimes\mathcal H_f$}
\label{sec:diag_offdiag}

The 9-dimensional past$\otimes$future Hilbert space carries a canonical
decomposition that incarnates the master equation $q! = 2q$ at the level of
quantum histories.

\begin{definition}[Diagonal and off-diagonal subspaces]\label{def:diag_subspaces}
For $\mathcal H_p\otimes\mathcal H_f = \CC^q\otimes\CC^q$, define
\begin{align}
  \mathcal D &\;=\; \mathrm{span}\bigl\{\ket{jj} : j\in\GF{q}\bigr\}
    & \text{(diagonal ``now'' subspace),}
  \\
  \mathcal N &\;=\; \mathrm{span}\bigl\{\ket{jk} : j,k\in\GF{q},\ j\neq k\bigr\}
    & \text{(off-diagonal history subspace).}
\end{align}
Then $\mathcal H_p\otimes\mathcal H_f = \mathcal D \oplus \mathcal N$ with
$\dim\mathcal D = q$ and $\dim\mathcal N = q^2 - q = q(q-1)$.
\end{definition}

\begin{theorem}[Master-equation history split]\label{thm:masteq_history}
For $q = 3$, the past/future Hilbert space splits as
\begin{equation}
  q^2  \;=\;  q + q!,
  \qquad
  9  \;=\;  3 + 6,
  \label{eq:masteq_history}
\end{equation}
where $3 = q = \dim\mathcal D$ is the diagonal ``now'' subspace and
$6 = q! = \dim\mathcal N$ is the off-diagonal history subspace.  This identity
holds \emph{only} at $q = 3$:
\begin{enumerate}[label=(\roman*),noitemsep]
  \item $\dim\mathcal D = q$ always; this is the simultaneous past$=$future
    eigenbasis, the ``simultaneous-identification'' subspace.
  \item $\dim\mathcal N = q(q-1)$ always; this is the
    ``context-changing'' subspace where past $\neq$ future.
  \item $q^2 = q + q(q-1)$ is trivial arithmetic, but
    $q(q-1) = q!$ is the \emph{master equation $q! = 2q$ rearranged}
    (multiplied by $(q-1)/(q-1)$), which holds uniquely at $q = 3$ in
    positive integers.
\end{enumerate}
\end{theorem}

\begin{proof}
(i)--(ii) by direct count.  (iii) is the rearrangement $q! = 2q \Leftrightarrow
q(q-1)(q-2)! = 2 \Leftrightarrow (q-2)! = 2/(q-1)$, which has integer solution
only at $q=3$ (giving $(q-2)! = 1! = 1 = 2/2$).  Equivalently, the master
equation \cref{thm:dioph} forces $q!=2q$, so for $q=3$ the off-diagonal count
$q(q-1) = q!\cdot(q-1)/q = q!$ when $(q-1) = q-1 = 2 = q!/q$, automatic at
$q=3$.
\end{proof}

\begin{corollary}[Bell qutrit is in the diagonal]\label{cor:bell_in_diagonal}
The Bell qutrit $\ket{\Omega} = q^{-1/2}\sum_j\ket{jj}$ lies entirely in
$\mathcal D$.  Its orthogonal complement in $\mathcal D$ (the other $q-1$
diagonal states) is the ``coherent diagonal cloud'' that carries the rest of the
Bell line's stabiliser action.
\end{corollary}

\begin{remark}[Three equivalent readings of $6$]
The off-diagonal dimension $6 = q!$ admits three substrate readings:
\begin{enumerate}[label=(\alph*),noitemsep]
  \item $q! = 6$ permutations of $\{0,1,2\}$ (symmetric group $S_3$ on the
    qutrit basis);
  \item $2q = 6$ ordered pairs (past, future) with past $\ne$ future (since
    each diagonal removes 1 of $q$ values, leaving $q-1$ futures for each
    past, giving $q(q-1) = 2q$ off-diagonal pairs at $q=3$);
  \item $q(q-1) = 6$ unordered choices of (past$\ne$future) pairs counted twice
    (once for each direction).
\end{enumerate}
These three readings of $6$ coincide \emph{only} at $q = 3$ via the master
equation, making the dimension of the off-diagonal history subspace a
substrate-primitive quantity.
\end{remark}

\subsection{One Photon vs Two Photons: Why Self-Entanglement Suffices}
\label{sec:one_vs_two}

The Bell state used in standard quantum optics is a \emph{two-photon} entangled
pair, $\ket{\Phi}_{AB} = q^{-1/2}\sum_j\ket{j}_A\otimes\ket{j}_B$, where photons
$A$ and $B$ are spatially separated.  The self-entangled state
\eqref{eq:bellqutrit} is structurally identical---the same Schmidt coefficients,
the same stabiliser group---but is implemented on \emph{one} photon, between
its past and future temporal-mode registers.

\begin{theorem}[Single-photon completeness]\label{thm:1photon_completeness}
Let a photon carry $r$ independent (mutually commuting) qutrit-capable degrees
of freedom---e.g.\ polarisation extended to qutrit, three-path spatial mode,
three time-bin temporal mode, three-sideband spectral mode, OAM ladder
$\ell\in\{-1,0,+1\}$.  Its mode Hilbert space then contains
$\CC^q \otimes \cdots \otimes \CC^q = \CC^{q^r}$ as a subspace.
For $r \ge 2$, the two-qutrit Pauli phase space $(\GF{q})^{2r}/\sim$ acting on
$\CC^{q^r}$ is, at $r = 2$ and $q = 3$, exactly the $W(3,3) = \Sp(4,\GF{3})$
substrate.  Consequently \emph{two independent qutrit registers on one photon
are sufficient to host the full $W(3,3)$ phase space}, and the past/future
Bell qutrit of \eqref{eq:bellqutrit} on one of those registers stabilises a
line of $W(3,3)$ exactly as in the two-photon case.
\end{theorem}

\begin{proof}
For $r = 2$ the symplectic Pauli phase space on $\CC^q\otimes\CC^q$ is
$(\GF{q})^4 / (\text{centre})$, and its non-zero rays modulo scalars form
$(3^4 - 1)/(3-1) = 40$ points with $40$ totally isotropic lines under the
canonical alternating form---this is the polar space $W(3,3)$
(\cref{prop:pts}).  The Bell qutrit's stabiliser is a totally isotropic
2-space (\cref{thm:bellchoi}), and the line/spread combinatorics are
identical regardless of whether the two qutrit registers are carried by one
photon or two.  Hence the self-entanglement and the two-photon entanglement
are operationally indistinguishable at the level of $W(3,3)$.
\end{proof}

\begin{corollary}[Photonic mode-DOF excess at $r \ge 2$]\label{cor:dof_excess}
A typical photon admits $r \ge 4$ commuting qutrit-capable registers:
polarisation$+$path (after qutrit-lifting via the vacuum mode), three time-bins,
three sidebands, OAM ladder.  Any \emph{two} of these suffice for $W(3,3)$;
the remaining $r - 2 \ge 2$ registers form a \emph{redundant} bookkeeping
shell, available for error correction and protected routing.  This excess is
the structural reason single-photon LOQC is fault-tolerant: the carrier is
already over-determined by two factors of $\CC^q = \CC^3$ beyond what the
W(3,3) substrate requires.
\end{corollary}

\begin{remark}[Operational meaning of past/future on one photon]
Past and future temporal modes on one photon are accessed
\emph{interferometrically}: the photon is split into two paths of unequal
length, then recombined.  The early-arriving amplitude is the ``past''
register; the late-arriving is the ``future'' register; the recombination
node is the ``now'' operation \eqref{eq:choi}, where the Bell-Choi overlap
$\mathrm{Tr}(U)/q$ is extracted as the interference visibility for any
inserted unitary $U$.  This is the standard Franson-interferometer geometry
\cite{Brendel1999,Marcikic2002}, and the substrate now provides the
group-theoretic interpretation: each Franson visibility measurement is a
sample of the Choi-Jamio\l kowski distribution of $\mathrm{U}(3)$.
\end{remark}

\subsection{Experimental Witness: Trace-Choi Visibility and Witting Contextuality}
\label{sec:witness}

The structure of \cref{thm:bell_uniqueness,thm:masteq_history,thm:1photon_completeness}
makes specific, falsifiable experimental predictions for any photonic
implementation of the temporal Bell qutrit.  We collect two independent
witnesses here.

\paragraph{Witness 1: Trace-Choi visibility.}
Insert a unitary $U \in \mathrm{U}(3)$ on the future-mode register of a
photon prepared in the temporal Bell qutrit $\ket{\Omega}$, then recombine
past and future at the Franson interferometer node.  The interference
visibility $V(U)$ between the diagonal and off-diagonal output ports is, by
\eqref{eq:choi},
\begin{equation}
  V(U) \;=\; \frac{|\mathrm{Tr}(U)|}{q} \;=\; \frac{|\mathrm{Tr}(U)|}{3}.
  \label{eq:trace_visibility}
\end{equation}

\begin{theorem}[Trace-Choi witness]\label{thm:trace_witness}
A photon source emits the temporal Bell qutrit $\ket{\Omega}$ \emph{iff} for
all $U \in \mathrm{U}(3)$ the Franson-output visibility satisfies
\eqref{eq:trace_visibility}.  In particular:
\begin{itemize}[noitemsep]
  \item $U = I$: $V = 1$ (maximal interference).
  \item $U = e^{i\alpha}I$: $V = 1$ (global phase preserves visibility).
  \item $U = X$ (cyclic shift $\ket{j}\mapsto\ket{j+1}$): $\mathrm{Tr}(X) = 0$
    so $V = 0$ (extinction).
  \item $U = Z$ (clock $\ket{j}\mapsto\omega^j\ket{j}$):
    $\mathrm{Tr}(Z) = 1+\omega+\omega^2 = 0$, $V = 0$.
  \item $U = F_3$ (qutrit Fourier): $|\mathrm{Tr}(F_3)|/3$ gives a substrate-
    specific intermediate visibility, computed in \cref{prop:F3vis}.
\end{itemize}
Any deviation from \eqref{eq:trace_visibility} certifies mixed-state
corruption or decoherence of the temporal entanglement.
\end{theorem}

\begin{proposition}[Fourier-gate visibility]\label{prop:F3vis}
For the qutrit Fourier gate $F_3$ with $F_3\ket{j} = q^{-1/2}\sum_k\omega^{jk}\ket{k}$,
\[
  \mathrm{Tr}(F_3) \;=\; \sum_{j=0}^{2} \frac{\omega^{j^2}}{\sqrt 3}
  \;=\; \frac{1 + \omega + \omega^4}{\sqrt 3}
  \;=\; \frac{1 + 2\omega}{\sqrt 3}
\]
since $\omega^4 = \omega$.  Hence
$|\mathrm{Tr}(F_3)| = |1+2\omega|/\sqrt 3 = \sqrt{3}/\sqrt 3 = 1$, giving
\begin{equation}
  V(F_3) \;=\; \frac{1}{3} \;=\; \frac{1}{q}.
\end{equation}
This is the \emph{quadratic Gauss sum} at $q=3$, and it provides the
characteristic substrate visibility for any inserted qutrit Fourier.
\end{proposition}

\paragraph{Witness 2: Witting non-contextuality.}
The 40 points of $W(3,3)$ embed canonically as the 40 \emph{Witting rays}
in $\CC^4$ via the (Vlasov 2025) symplectic realisation; any pair
$\ket{\psi},\ket{\varphi}$ of Witting rays satisfies the discrete inner-product
condition
\begin{equation}
  |\langle\psi | \varphi\rangle|^2 \;\in\; \{0,\ \tfrac{1}{3}\}
  \;=\; \{0,\ 1/q\}.
  \label{eq:witting_overlap}
\end{equation}
The orthogonal pairs ($|\langle\psi|\varphi\rangle|^2 = 0$) correspond to
collinear Pauli observables on the $W(3,3)$ line structure; the
$1/q$-overlap pairs correspond to non-collinear observables.

\begin{theorem}[Witting Kochen--Specker bound]\label{thm:KSbound}
No assignment of definite $\{0,1\}$-values to the $40$ Witting rays can be
consistent with the orthogonality structure: every $\{0,1\}$-marking is
\emph{Kochen--Specker contradictory} on at least $6 = q!$ tetrads.  The
maximum number of \emph{satisfiable} mutually-exclusive Witting tetrads is
$34/40$, and the marking optima form an orbit of size $720 = (q!)!$
(Vlasov 2025).
\end{theorem}

In a single-photon experiment with the temporal Bell qutrit, a sequence of
$N$ Pauli measurements drawn from a Witting tetrad covering of $W(3,3)$
will exhibit \emph{exactly} $\le 34N/40$ classically-consistent assignments;
any larger fraction certifies contextual (genuinely quantum) substrate
content.  Conversely, $\ge 35N/40$ classical agreement falsifies the
W(3,3) self-entanglement claim.

\paragraph{Witness 3: Key-agreement rate.}
A two-party photonic protocol using $W(3,3)$ as the shared Pauli phase space
has a unique key-agreement probability per round:
\begin{equation}
  P_{\rm key}(W(3,3)) \;=\; \frac{13}{40} \;=\; \frac{\Phi_3}{v},
  \label{eq:keyrate}
\end{equation}
matching the Vlasov 2025 Table III count of $40 = 28 + 12 = \mu\Phi_6 + k$
basis decompositions.  A measured rate $13/40$ over many rounds is a third
falsifiable substrate signature distinct from the trace-Choi and Witting
witnesses.

These three witnesses---trace-Choi visibility, Witting non-contextuality,
and the $13/40$ key-agreement rate---are mutually independent
falsifiable predictions of the single-photon self-entanglement framework.

\subsection{Clifford Action: Bell-Line Orbit and Stabiliser}
\label{sec:clifford_orbit}

The two-qutrit Clifford group $C_2 = \Sp(4,\GF{3})$ (of order $51840$) acts
naturally on the 40 lines of $W(3,3)$.  The Bell-line stabiliser
$\langle X_p X_f,\ Z_p Z_f^{-1}\rangle$ identified in \cref{thm:bellchoi} sits
inside one specific orbit under this action; we compute the orbit and
stabiliser structure here.

\begin{theorem}[Bell-line orbit decomposition]\label{thm:orbit}
The two-qutrit Clifford group $\Sp(4,\GF{3})$ acts transitively on the 40
lines of $W(3,3)$, so the orbit of the Bell line is all of $W(3,3)$:
\begin{equation}
  |\mathrm{orbit}(\mathrm{Bell\ line})| \;=\; 40 \;=\; v.
\end{equation}
The stabiliser of the Bell line is the subgroup of $\Sp(4,\GF{3})$ fixing
the totally isotropic 2-space $\langle X_p X_f,\ Z_p Z_f^{-1}\rangle$, of order
\begin{equation}
  |\mathrm{Stab}(\mathrm{Bell\ line})|
  \;=\; \frac{|\Sp(4,\GF{3})|}{|\mathrm{orbit}|}
  \;=\; \frac{51840}{40}
  \;=\; 1296
  \;=\; \mu^2 \cdot q^{q+1}.
\end{equation}
That is, the Bell-line stabiliser has order equal to $\mu^2$ times the
matter-sector dimension $q^{q+1} = H_1(2\text{-complex})$.
\end{theorem}

\begin{proof}
Transitivity of $\Sp(4,\GF{3})$ on isotropic 2-spaces of $\GF{3}^4$ is
standard (the symplectic Witt extension theorem~\cite{Tits1974}).  The
orbit-stabiliser ratio is $|G|/|\mathrm{orbit}| = 51840/40 = 1296$.
Factoring: $1296 = 16 \cdot 81 = \mu^2 \cdot q^{q+1}$, where $\mu^2 = 16 = 2^\mu$
and $q^{q+1} = 3^4 = 81$ is the matter sector dimension (the substrate
$H_1$ of the W(3,3) line-triangle 2-complex, established via the Hodge
decomposition).
\end{proof}

\begin{corollary}[Substrate factorisation of the Clifford group]\label{cor:cliff_factor}
The two-qutrit Clifford group decomposes via the Bell-line action as
\begin{equation}
  |\Sp(4,\GF{3})| \;=\; |\mathrm{orbit}| \cdot |\mathrm{Stab}|
  \;=\; v \cdot \mu^2 \cdot q^{q+1}
  \;=\; 40 \cdot 16 \cdot 81.
\end{equation}
Three substrate primitives $\{v, \mu^2, q^{q+1}\}$ exactly multiplicatively
exhaust the Clifford group order.
\end{corollary}

\begin{remark}[Operational interpretation]
The Bell-line stabiliser of order $\mu^2 \cdot q^{q+1} = 1296$ is precisely
the group of two-qutrit Clifford gates that preserve the temporal Bell
qutrit's stabiliser line---i.e., the gates that respect the
\emph{past$\leftrightarrow$future symmetry} of the temporal seed.  Outside
the stabiliser, the remaining $51840 - 1296 = 50544$ Cliffords \emph{rotate
the Bell line through the W(3,3) line graph}, scrambling the temporal seed
across the 40-vertex substrate.  This $1296 : 50544 \approx 1 : 39$
ratio (matching $1 : (v-1)$) is the substrate measure of how much
``temporal symmetry preservation'' costs in the Clifford computational
budget.
\end{remark}

\subsection{Photonic Mode-Space Requirement: $q^4$ as the Substrate Threshold}
\label{sec:mode_space_req}

A single photon's mode space must be ``wide enough'' to host the W(3,3)
substrate.  We quantify this requirement in terms of standard photonic
phase-space variables (frequency $\omega$, time $t$, polarisation, OAM,
sideband).

\begin{theorem}[Mode-space threshold]\label{thm:mode_space}
A single photon supports the full W(3,3) self-entanglement compilation
of \cref{thm:selfent} \emph{iff} its mode-Hilbert-space dimension satisfies
\begin{equation}
  \dim \mathcal H_{\rm photon}  \;\ge\;  q^{r}
  \;\text{where }r \ge 2,
  \qquad
  \text{with }r = 4 \text{ realising the full }W(3,3)\text{ phase space.}
  \label{eq:mode_space_req}
\end{equation}
At the cleanest threshold $r = 4$:
\begin{equation}
  \dim \mathcal H_{\rm photon} \;\ge\; q^4 \;=\; 81 \;=\; q^{q+1}
  \;=\; H_1(2\text{-complex }W(3,3)),
\end{equation}
i.e., the photon's mode-space dimension threshold equals exactly the
W(3,3) matter-sector dimension.
\end{theorem}

\begin{proof}
By \cref{thm:1photon_completeness}, $r = 2$ commuting qutrit registers
suffice for $W(3,3)$.  For the full $\GF{3}^4$ Pauli phase space lifted to a
mode Hilbert space (not modulo scalars), we need $r = 4$ commuting
qutrit-capable DOFs, giving $\dim = q^4 = 81$.  And $q^4 = q^{q+1} = 81 =
H_1(W(3,3)\ 2\text{-complex})$ is the W(3,3) matter sector
(commit \texttt{ac4dfadc}).
\end{proof}

\begin{corollary}[Time-bandwidth product]\label{cor:tb_product}
For a photon with $q = 3$ time-bins of width $\Delta t_{\rm bin}$ and
$q = 3$ spectral sidebands of width $\Delta\omega_{\rm sideband}$
(temporal $\otimes$ spectral two-qutrit register), the photonic
time-bandwidth product must satisfy
\begin{equation}
  \Delta\omega_{\rm total} \cdot \Delta t_{\rm total}
  \;\gtrsim\;
  q^2 \;=\; 9.
\end{equation}
For the full $r = 4$ realisation (polarisation $\times$ path $\times$
time-bin $\times$ sideband):
\begin{equation}
  \Delta\omega \cdot \Delta t \;\gtrsim\; q^4 \;=\; 81.
\end{equation}
\end{corollary}

\begin{remark}[Substrate equates matter sector with photonic envelope]
The threshold $q^4 = 81$ is the W(3,3) matter sector dimension (the
harmonic 1-form count, equivalently the dim of $H_1$ of the line-triangle
2-complex).  Hence the photonic substrate \emph{matches its matter sector
to its envelope phase space}: the smallest single-photon envelope capable
of hosting the W(3,3) substrate has time-bandwidth product equal to the
substrate's matter-sector count.  This identifies the photonic substrate
threshold as a topological quantity ($H_1$) rather than a free
phenomenological parameter.
\end{remark}

\begin{remark}[Wavelength domain estimate]
For typical telecom-band single photons at $\lambda = 1550$ nm
($\omega_0 \sim 1.2 \times 10^{15}$ rad/s) with $\Delta t \sim 1$ ns
time-bin spacing, the time-bandwidth product is
$\Delta\omega \cdot \Delta t \sim 10^{6}$, vastly exceeding the W(3,3)
threshold $q^4 = 81$.  Femtosecond-scale photons reduce this by a factor
$10^3$ to $\sim 10^3$, still well above threshold.  Hence \emph{all
practically generated single photons satisfy the W(3,3) mode-space
threshold by 1--3 orders of magnitude}.
\end{remark}

\subsection{Quantum Teleportation in Time}
\label{sec:time_teleportation}

The most operational use of the temporal Bell qutrit is to transmit a qutrit
state from a past-self to a future-self on a single photon---a temporal
analog of standard quantum teleportation.

\begin{theorem}[Single-photon temporal teleportation]\label{thm:temporal_teleport}
Let a single photon carry the temporal Bell qutrit $\ket{\Omega}$ on
its past$\otimes$future registers, and let $\ket{\psi} \in \CC^q$ be an
arbitrary qutrit state encoded on an auxiliary qutrit register.  Then
$\ket{\psi}$ can be teleported from the past-self's accessible region to
the future-self's accessible region with fidelity $1$, via the protocol:
\begin{enumerate}[label=(\roman*),noitemsep]
  \item Past-self performs a qutrit Bell-basis measurement on
    $\ket{\psi}_{\rm aux}\otimes\ket{\Omega_p}$ (the auxiliary qutrit jointly
    with the past register of the Bell qutrit).
  \item Past-self stores the $2$-trit measurement outcome in classical
    memory.
  \item At a later time, future-self reads the classical memory and applies
    the corresponding Pauli correction
    $X^a Z^b$ to the future register of the Bell qutrit, where $(a,b)$ is the
    measurement outcome.
  \item The future register now carries $\ket{\psi}$ exactly.
\end{enumerate}
The protocol consumes one temporal Bell qutrit and $2 \log_q q = 2$ classical
trits ($= 2 \log_2 3 \approx 3.17$ bits) of memory.
\end{theorem}

\begin{proof}
Standard quantum teleportation theorem~\cite{Bennett1993}, adapted to qutrits
and to the temporal-Bell setting.  The classical channel from past-self to
future-self is realised by classical memory (records persisting in time)
rather than spatial communication; the rest of the protocol is identical.
\end{proof}

\begin{corollary}[Single-photon memory protocol]\label{cor:memory_protocol}
A photon source preparing temporal Bell qutrits, together with a Bell-basis
detector, a $2$-trit classical memory, and a qutrit Pauli-correction unit,
implements a faithful qutrit quantum memory: any qutrit state encoded at
time $t_0$ can be retrieved at time $t_1 > t_0$ on the same photon, with no
additional photons required.  The classical-memory overhead per stored qutrit
is $2$ trits, independent of storage time.
\end{corollary}

\begin{remark}[No retrocausation]
The protocol does \emph{not} require backwards-in-time signalling: classical
information flows forward (past stores, future reads), and the quantum state
is transmitted via the pre-existing past$\otimes$future entanglement.  This
is consistent with standard quantum-information no-signalling.  The temporal
Bell qutrit's ``self-entanglement'' is a static structural property of the
photon's temporal-mode register, not a dynamical signal.
\end{remark}

\subsection{Decoherence Threshold: Substrate-Clean Noise Tolerance}
\label{sec:decoherence}

A realistic photonic implementation suffers depolarising noise on the
temporal qutrit registers.  The Bell qutrit's entanglement survives noise
up to a substrate-clean threshold.

\begin{theorem}[Bell-qutrit decoherence threshold]\label{thm:noise_threshold}
Consider the noisy temporal Bell state (isotropic Werner-like form)
\begin{equation}
  \rho(p) \;=\; (1-p)\,\ket{\Omega}\bra{\Omega} + \frac{p}{q^2}\, I_{q^2},
  \qquad p \in [0, 1].
  \label{eq:werner}
\end{equation}
Then $\rho(p)$ is past$\leftrightarrow$future entangled \emph{if and only if}
\begin{equation}
  p \;<\; \frac{q}{q+1} \;=\; \frac{q}{\mu} \;=\; \frac{3}{4}.
  \label{eq:noise_thresh}
\end{equation}
The substrate noise threshold is exactly $q/\mu = 3/4$.
\end{theorem}

\begin{proof}
Standard Werner-Holevo result: an isotropic two-qudit state
$\rho(\lambda) = \lambda\ket{\Omega}\bra{\Omega} + (1-\lambda)/q^2 \cdot I$
is separable iff $\lambda \le 1/(q+1)$.  Reparametrising via $\lambda = 1-p$,
separability holds iff $1-p \le 1/(q+1)$, i.e.\ $p \ge q/(q+1) = q/\mu$.
At $q = 3$: $p_{\text{sep}} = 3/4$.
\end{proof}

\begin{corollary}[Substrate noise readings]\label{cor:noise_readings}
The threshold $3/4 = q/\mu$ admits four substrate-primitive readings:
\begin{enumerate}[label=(\roman*),noitemsep]
  \item $q/\mu$ (fundamental quantum / co-quantum),
  \item $(q+1-1)/(q+1) = 1 - 1/\mu$ (one minus inverse co-quantum),
  \item $(\lambda+1)/\mu = (\text{SRG }\lambda + 1)/\mu$ (SRG $\lambda$ shifted by one),
  \item $(\mu-1)/\mu$ (co-quantum minus one over co-quantum).
\end{enumerate}
\end{corollary}

\begin{remark}[Comparison with two-qubit noise threshold]\label{rem:qubit_compare}
For two-qubit ($q=2$) Werner states, the separability threshold is
$p = q/(q+1) = 2/3$ via the same formula.  The W(3,3)-qutrit threshold
$3/4 > 2/3$ makes the temporal Bell qutrit \emph{strictly more
noise-tolerant} than the standard qubit Bell pair: the substrate jump from
$q=2$ to $q=3$ buys noise tolerance from $2/3$ to $3/4$, a gain consistent
with the substrate's master equation $q! = 2q$ activating at $q=3$.
\end{remark}

\begin{remark}[Quantum capacity below threshold]\label{rem:capacity}
For $p < q/\mu = 3/4$, the entanglement of formation of $\rho(p)$ is positive
and the joint state supports quantum teleportation, dense coding, and other
entanglement-based protocols.  At $p = q/\mu$, entanglement degrades to zero
and the joint state becomes a separable mixture of product states.
\end{remark}

\subsection{Marginal Mixedness and Mutual Information}
\label{sec:marginal}

The temporal Bell qutrit hides all its information in the correlations between
past and future, leaving the individual marginal states maximally mixed.

\begin{theorem}[Marginal mixedness]\label{thm:marginal}
For the temporal Bell qutrit $\ket{\Omega} = q^{-1/2}\sum_j\ket{jj}$:
\begin{enumerate}[label=(\roman*),noitemsep]
  \item $\rho_p = \mathrm{Tr}_f\bigl(\ket{\Omega}\bra{\Omega}\bigr)
    = \frac{1}{q}\,I_q$, the maximally mixed qutrit state on the past register;
  \item $\rho_f = \mathrm{Tr}_p\bigl(\ket{\Omega}\bra{\Omega}\bigr) = \frac{1}{q}\,I_q$
    on the future register;
  \item von Neumann entropy $S(\rho_p) = S(\rho_f) = \log q = \log 3$;
  \item entropy of the joint state $S(\ket{\Omega}\bra{\Omega}) = 0$ (pure);
  \item quantum mutual information $I(p:f) = 2 \log q = 2 \log 3$;
  \item coherent information $I_c(p \rightarrow f) = \log q$ (entanglement
    entropy).
\end{enumerate}
\end{theorem}

\begin{proof}
$\rho_p = \mathrm{Tr}_f(\frac{1}{q}\sum_{j,k}\ket{jj}\bra{kk})
= \frac{1}{q}\sum_{j,k}\delta_{jk}\ket{j}\bra{k}
= \frac{1}{q}\sum_j\ket{j}\bra{j} = I_q/q$.  Similarly $\rho_f = I_q/q$.
The von Neumann entropy of the maximally mixed qutrit is $\log q$, and
$S(\ket{\Omega}) = 0$ since $\ket{\Omega}$ is pure.  Mutual information
$I(p:f) = S(\rho_p) + S(\rho_f) - S(\ket{\Omega}\bra{\Omega}) = 2 \log q$,
and coherent information for pure states equals $S(\rho_p)$.
\end{proof}

\begin{corollary}[All information is in correlations]\label{cor:correl}
A single-time measurement of either the past or the future register alone
returns a uniformly random outcome in $\GF{q}$.  The temporal Bell qutrit's
\emph{entire} information content lives in the past-future correlation; no
single-time observer can recover it.
\end{corollary}

\begin{remark}[Substrate principle: ``past alone is noise'']
The marginal-mixedness theorem encodes a deep substrate principle: at the
single-photon temporal level, past-alone-measurements and
future-alone-measurements are individually \emph{maximally noisy}.  The
information is reconstructed only by measuring the joint past$\otimes$future
state---equivalently, by the Franson-recombination geometry of
\cref{rem:phys_sym}.  This is the operational meaning of ``now is the
harmonic convergence'': only at the convergence point (where past and
future are read together) does the substrate's information become accessible.
\end{remark}

\begin{remark}[Mutual information bound: $I(p:f) = 2 \log q$]
For any qutrit-qutrit state, the quantum mutual information $I(p:f)$ is
bounded above by $2 \log q$; the temporal Bell qutrit saturates this bound.
This makes $\ket{\Omega}$ \emph{the unique state} (up to local unitaries)
that achieves maximum past$\leftrightarrow$future mutual information, an
independent characterisation complementing \cref{thm:bell_uniqueness}.
\end{remark}

\subsection{Preparation Protocol: From $\ket{00}$ to $\ket{\Omega}$ in Two Cliffords}
\label{sec:prep_protocol}

The temporal Bell qutrit \eqref{eq:bellqutrit} is prepared from the trivial
product state $\ket{0}_p \ket{0}_f$ by a depth-2 two-qutrit Clifford circuit
consisting of one qutrit Hadamard $F_3$ on the past register followed by a
controlled cyclic shift between past and future.

\begin{proposition}[Bell qutrit preparation]\label{prop:prep}
Define the qutrit Hadamard and controlled-shift gates by
\begin{align}
  F_3\ket{j} &\;=\; q^{-1/2}\sum_k \omega^{jk}\ket{k},
  &\omega &= e^{2\pi i / q},
  \\
  CX_{p\to f}\ket{j}_p\ket{k}_f &\;=\; \ket{j}_p\ket{k + j \bmod q}_f.
\end{align}
Then
\begin{equation}
  CX_{p\to f} \,\bigl(F_3 \otimes I\bigr)\,\ket{0}_p\ket{0}_f
  \;=\; \frac{1}{\sqrt q}\sum_{j=0}^{q-1}\ket{j}_p\ket{j}_f
  \;=\; \ket{\Omega}.
  \label{eq:prep_recipe}
\end{equation}
\end{proposition}

\begin{proof}
Direct computation: $F_3\ket{0}_p\ket{0}_f =
q^{-1/2}\sum_j\omega^{0\cdot j}\ket{j}_p\ket{0}_f = q^{-1/2}\sum_j\ket{j}_p\ket{0}_f$.
Then $CX_{p\to f}$ takes $\ket{j}_p\ket{0}_f \mapsto \ket{j}_p\ket{j}_f$, giving
\eqref{eq:prep_recipe}.
\end{proof}

\begin{remark}[Photonic time-bin implementation]\label{rem:photonic_prep}
For a single photon with three time-bin modes:
\begin{enumerate}[label=(\alph*),noitemsep]
  \item Initial state: photon in time-bin $0$ on both past and future
    registers (e.g.\ via electro-optic preselection).
  \item Past Hadamard ($F_3$): pass the past register through a 3-port
    tritter (e.g.\ a balanced fiber-optic tritter or a three-prism beam
    splitter), creating a uniform superposition of past-bins.
  \item Controlled-shift ($CX_{p\to f}$): use a phase modulator triggered by
    the past time-bin to cyclically permute the future time-bin (e.g.\ via
    electro-optic modulation conditioned on detection of past-bin tag).
  \item Output state: the temporal Bell qutrit $\ket{\Omega}$.
\end{enumerate}
The preparation is fully linear-optical (no measurement-induced
nonlinearity), requires only one tritter and one conditional phase modulator,
and is depth-$q$ in time-bin operations.
\end{remark}

\begin{corollary}[Two-Clifford depth complexity]\label{cor:depth2}
The temporal Bell qutrit is prepared in \emph{exactly two} Clifford
operations from the trivial state.  This matches the qutrit-Bell-state
preparation complexity ($F + CX$) and is the minimum for a maximally
entangled state on $\CC^q \otimes \CC^q$: any preparation requires at least
$O(\log q)$ depth on a fully connected register, achieved here at $O(1)$
because $F_3$ creates the diagonal cloud and $CX$ folds it into the Bell
basis.
\end{corollary}

\subsection{The Triple Closure: $7$, $9$, $40$ on One Photon}
\label{sec:triple_closure}

The single-photon self-entanglement compiles to the full $W(3,3)$ substrate
through three layered closures, each carrying a different substrate primitive.

\begin{theorem}[Temporal triple closure]\label{thm:triple_closure}
At $q = 3$, a single self-entangled qutrit on one photon realises three nested
substrate closures simultaneously:
\begin{enumerate}[label=(\arabic*),noitemsep]
  \item \textbf{Temporal triangle ($\Phi_6 = 7$):}
    \begin{equation}
      \underbrace{3}_{\text{past cells}} + \underbrace{3}_{\text{future cells}}
      + \underbrace{1}_{\text{now}} \;=\; 7 \;=\; \Phi_6.
    \end{equation}
  \item \textbf{Past$\otimes$future Hilbert space ($q^2 = 9$):}
    \begin{equation}
      q + q!  \;=\; 3 + 6 \;=\; 9 \;=\; q^2.
      \qquad (\text{diagonal} + \text{off-diagonal})
    \end{equation}
  \item \textbf{$W(3,3)$ shell ($v = 40$):}
    \begin{equation}
      1 + 12 + 27 \;=\; 40 \;=\; v,
      \qquad (\text{Bell line} + \text{intersecting} + \text{disjoint})
    \end{equation}
    with substrate readings $12 = k$ (valency) and $27 = q^q$ (Heisenberg-Weyl
    order).
\end{enumerate}
The three closures are nested via:
\begin{equation}
  \mathrm{temporal\ triangle}_{(7)}  \;\hookrightarrow\;
  \mathrm{past}\otimes\mathrm{future}_{(9)}  \;\hookrightarrow\;
  W(3,3)_{(40)}.
\end{equation}
\end{theorem}

\begin{proof}
(1) The temporal triangle decomposes a single self-entangled qutrit into $q$
past cells (one per substrate value), $q$ future cells, and a single
``now'' point---the harmonic convergence of \cref{cor:now_forced}.  Total:
$2q + 1 = 7 = \Phi_6$.
(2) Established in \cref{thm:masteq_history}.
(3) The Bell line is one of $v = 40$ lines; the Bell-local shell closes as
$1 + 12 + 27 = 40$ by \cref{thm:nowclosure}.  Substrate readings:
$12 = k$ is the W(3,3) valency, $27 = q^q$ is the Heisenberg-Weyl group
order.
The nesting follows from the embeddings $\mathrm{cells} \to
\mathrm{Hilbert\ space} \to \mathrm{phase\ space}$: the $7$ temporal cells
generate the $9$-dimensional past$\otimes$future Hilbert space, which then
acts via Pauli observables on the $40$-point $W(3,3)$ phase space.
\end{proof}

\begin{corollary}[Total substrate scale-up]\label{cor:scaleup}
Starting from $1$ photon with $1$ temporal qutrit register ($q = 3$ time-bins),
the substrate scales up to:
\begin{center}\begin{tabular}{lll}
$1$ photon &           & seed \\
$q = 3$    & temporal modes & past, now, future cells \\
$q^2 = 9$  & past$\otimes$future histories & diagonal $+$ off-diagonal \\
$v = 40$   & W(3,3) phase-space points & full substrate \\
$q^{q+1} = 81$ & matter sector & harmonic 1-forms \\
$|\Sp(4,\GF{3})| = 51840$ & two-qutrit Clifford & all gates \\
\end{tabular}\end{center}
Each step is a substrate-clean multiplicative or combinatorial scale,
ultimately exhausting the full universal-quantum-computation substrate from
a single photon's temporal qutrit.
\end{corollary}

\subsection{Substrate CPT Theorem on the Bell Qutrit}
\label{sec:cpt}

The temporal Bell qutrit is invariant under a discrete CPT-like
symmetry, which combines complex conjugation, past$\leftrightarrow$future
parity exchange, and unitary time reversal.

\begin{theorem}[Substrate discrete CPT]\label{thm:cpt}
Define the three discrete symmetries on $\mathcal H_p\otimes\mathcal H_f$:
\begin{itemize}[noitemsep]
  \item $C$ (charge conjugation):  $\ket{j} \mapsto \ket{(-j)\bmod q}$ on
    each register, equivalently $\omega \mapsto \overline\omega$.
  \item $P$ (parity / past$\leftrightarrow$future):
    $\mathrm{SWAP}_{pf}$, exchanging the two registers.
  \item $T$ (time reversal):  complex conjugation
    $\ket{\psi} \mapsto \overline{\ket{\psi}}$ on coefficients (with $T^2 = I$
    on a qutrit-tensor system).
\end{itemize}
Then the temporal Bell qutrit is preserved by each of $C$, $P$, $T$
individually, and \emph{a fortiori} by the combined $CPT$:
\begin{equation}
  C\ket{\Omega} = P\ket{\Omega} = T\ket{\Omega} = CPT\ket{\Omega} = \ket{\Omega}.
\end{equation}
\end{theorem}

\begin{proof}
$C\ket{\Omega} = q^{-1/2}\sum_j\ket{-j,-j} = q^{-1/2}\sum_j\ket{j,j} =
\ket{\Omega}$ by reindexing $j' = -j$.
$P\ket{\Omega} = q^{-1/2}\sum_j\ket{j,j} = \ket{\Omega}$ since $\ket{j,j}$
is SWAP-symmetric.
$T\ket{\Omega} = q^{-1/2}\overline{\sum_j\ket{j,j}} = \ket{\Omega}$ since
$\ket{j,j}$ has real computational-basis coefficients.
Hence $CPT$ also fixes $\ket{\Omega}$.
\end{proof}

\begin{corollary}[$(U\otimes U^*)$ invariance as $CPT$]\label{cor:uust_cpt}
The unitary-conjugate invariance $(U\otimes U^*)\ket{\Omega} = \ket{\Omega}$
of \cref{thm:bell_uniqueness}(iii) is equivalent to the statement:
\emph{any forward evolution $U$ on the future register is matched by the
$C$-conjugate (= time-reversed) evolution $U^*$ on the past register.}
In particular, the temporal Bell qutrit cannot distinguish
``future evolving forward'' from ``past evolving backward''---the
substrate is $CPT$-symmetric at the single-photon level.
\end{corollary}

\begin{remark}[Lüders--Pauli analog]
\cref{thm:cpt} is the substrate's discrete-field analog of the
Lüders--Pauli CPT theorem of relativistic QFT: any consistent
self-entangled temporal state on a single photon is automatically
$CPT$-invariant.  This is forced by the (U$\otimes$U$^*$) characterisation
(\cref{thm:bell_uniqueness}) and does not require any continuum spacetime
input.  The temporal Bell qutrit thus realises a \emph{discrete-substrate
CPT symmetry} purely on photonic mode space.
\end{remark}

\subsection{Synthesis}

This establishes the physical thesis: \emph{one self-entangled qutrit on one single
photon} carries the full $W(3,3)$ substrate as its now-context spread, with the
Bell stabiliser line as the temporal seed.  All structural counts of
\cref{sec:w33}---$v=40$, $\lambda=2$, $\mu=4$, $q+1=4$ rays per line, $36$ complete
programs---read off directly from this self-entangled photonic carrier.

%======================================================================
\section{Universal Gate Sets}
\label{sec:gates}
%======================================================================

\subsection{Qubit Gates via Linear Optics}

A linear optical network acts on mode operators by a unitary $U \in \mathrm{U}(m)$:
$\hat{a}_k \mapsto \sum_l U_{kl}\,\hat{a}_l$.
For the polarisation qubit, the two generators are the wave plate (retardance
$\delta$, fast axis angle $\theta$)~\cite{Reck1994} and the balanced beam splitter
$U_\mathrm{BS} = \frac{1}{\sqrt{2}}\bigl(\begin{smallmatrix}1 & i \\ i & 1\end{smallmatrix}\bigr)$.
A half-wave plate and quarter-wave plate together generate all of $\mathrm{SU}(2)$.
Any single-qubit unitary is realisable with at most $q = 3$ wave plates---a direct
consequence of the $W(3,3)$ Diophantine parameter.

\subsection{Qutrit Gates: The Clifford Group \texorpdfstring{$\Sp(4,\GF{3})$}{Sp(4,F3)}}

For the photonic qutrit, the universal gate set is:
\begin{itemize}[noitemsep]
  \item Qutrit Hadamard $F_3$: quantum Fourier transform,
    $F_3\ket{j} = 3^{-1/2}\sum_k \omega^{jk}\ket{k}$.
  \item Controlled-$Z$ (qutrit): $CZ_3\ket{j,k} = \omega^{jk}\ket{j,k}$.
  \item Phase gate $S_3 = \mathrm{diag}(1, 1, \omega)$.
  \item Magic state injection via cubic phase gate
    $T_3 = \mathrm{diag}(1, \omega, \omega^4)$.
\end{itemize}
The group $\langle F_3, CZ_3, S_3 \rangle$ is exactly $\Sp(4,\FF_3)$; adding $T_3$
gives a universal gate set (Bravyi--Kitaev~\cite{BK2005}).

%======================================================================
\section{Non-Deterministic Two-Qubit Gates: The KLM Protocol}
\label{sec:klm}
%======================================================================

Single-photon states of fixed photon number evolve under mode-preserving linear
optics; because no photon--photon interaction exists in vacuum, a deterministic
controlled-NOT is impossible without additional resources.  Knill, Laflamme, and
Milburn~\cite{KLM2001} resolved this by demonstrating:
\begin{enumerate}[noitemsep]
  \item A \emph{probabilistic} CZ gate is implementable with ancilla vacuum modes
    and single-photon detection.
  \item The success probability approaches unity via \emph{teleportation} with
    offline resource states.
\end{enumerate}
Post-selection onto the $\ket{1}$ ancilla outcome realises
$\ket{1,1} \mapsto -\ket{1,1}$
at probability $1/4$ for the minimal circuit.  The number of ancilla modes in the
elementary KLM gadget is $\mu = 4$---matching the $\mu=4$ parameter of $W(3,3)$.

%======================================================================
\section{Measurement-Based Quantum Computation with Single Photons}
\label{sec:mbqc}
%======================================================================

Measurement-based quantum computing (MBQC)~\cite{Raussendorf2001} reframes
universality in terms of a pre-entangled \emph{cluster state} and adaptive
single-qubit measurements.  A cluster state $\ket{G}$ on graph $G=(V,E)$ satisfies
$K_a\ket{G} = \ket{G}$ for all $a\in V$, where $K_a = X_a \bigotimes_{b\sim a} Z_b$.

\textbf{Photon-fusion networks.}  Type-II fusion gates generate graph-state edges
between photonic qubits with probability $1/2$.  A 2D cluster state can be built
reliably whenever fusion success exceeds the bond-percolation threshold
$p_c \approx 0.5$~\cite{Browne2005}.

\textbf{The $W(3,3)$ cluster state.}  The collinearity graph of $W(3,3)$---40
vertices, 240 edges, diameter 2, valency 12---defines a natural photonic cluster
substrate: 40 photonic qubits (or qutrits) in dual-rail (or three-mode) encoding,
with entangling fusion gates applied along the 240 graph edges.  The Laplacian
spectral gap $\Delta_L = k - r = 12 - 2 = 10$ bounds the mixing time and information
propagation rate.  Universal MBQC on the $W(3,3)$ cluster state requires $v = 40$
photons and $E = 240$ entangling operations.

%======================================================================
\section{The Four-Layer Photonic Runtime}
\label{sec:runtime}
%======================================================================

The previous sections describe the physical ingredients separately: coherent
qutrits, probabilistic optical gates, and one-way measurement.  The architecture is
cleanest when these are kept as four different runtime layers.

\begin{theorem}[Photonic runtime separation]
\label{thm:runtime}
The $W(3,3)$ single-photon architecture separates into a quantum carrier, a
probabilistic assembly layer, a deterministic feed-forward layer, and a classical
record layer:
\[
\begin{array}{ccl}
\text{quantum carrier} &:& \GF{3}^{4}/\GF{3}^{\times},\quad
  (3^4-1)/(3-1)=40,\\
\text{probabilistic assembly} &:& p_{\mathrm{fusion}}=\lambda/\mu=1/2,\quad
  p_{\mathrm{KLM}}=1/\mu=1/4,\\
\text{deterministic logic} &:& \text{MBQC Pauli-frame correction on }3^4=81
  \text{ frame states},\\
\text{classical record} &:& 40\text{ trits with }2^{63}<3^{40}<2^{64}.
\end{array}
\]
\end{theorem}

\begin{proof}
The quantum carrier is \cref{thm:phasespace}: two-qutrit Pauli exponent vectors in
$\GF{3}^{4}$ projectivise to the 40 observables of $W(3,3)$.  Type-II fusion
succeeds with probability $1/2=\lambda/\mu$, so assembling all $E=240$ graph edges
requires expected budget $E/(1/2)=480=2E$, the directed Hashimoto carrier.  The
minimal KLM controlled-phase success probability is $1/4=1/\mu$, so the same edge
set has primitive KLM budget $4E=960=\mathrm{tr}(A^3)$.  MBQC measurements are
random at the detector, but deterministic at the logical layer because the
measurement word updates a Pauli frame; for two qutrits that frame has $3^4=81$
states, projectivising back to the 40-point carrier.  The complete single-run
measurement record is a 40-trit word, and direct calculation gives
$2^{63}<3^{40}<2^{64}$.
\end{proof}

\subsection{Topology and Feed-Forward}

The topological side of the same runtime is the genus-one toric/Csaszar--Szilassi
shell.  The toric code on a torus has $2g=2=\lambda$ logical qubits and ground-state
degeneracy $2^{2g}=4=\mu$.  The Csaszar $K_7$ triangulation has
$(V,E,F)=(7,21,14)$, Euler characteristic $0$, and genus $1$; its Szilassi dual
keeps the same toroidal code surface.  The Jungerman--Ringel complete-graph genus
denominator is $12=k$, and $K_{12}$ has orientable genus $6=k/\lambda$.

This is why the number $12$ in the photonic architecture should not be read only as
a Clifford/gauge count.  It is also the minimal-triangulation denominator governing
which complete graph surfaces can host the topological memory used by the
feed-forward layer.

\subsection{Hashimoto Phase Angles: Photonic Substrate Predictions}
\label{sec:hashimoto-phases}

The directed-edge Hashimoto operator $B$ from \cref{sec:hashimoto-phases}'s
context (assembling the $480$ directed carriers from the $E = 240$ undirected
edges) admits a sector-projected spectrum forced by the Ihara--Bass identity.
Each adjacency eigenvalue $\lambda$ of $W(3,3)$ contributes Hashimoto
eigenvalues from $u^2 - \lambda u + (k-1) = 0$; substituting the
$\SRG(40,12,2,4)$ spectrum $\{12, 2, -4\}$ gives the gauge-sector pair
\[
u_{\mathrm{gauge}} \;=\; 1 \;\pm\; i\sqrt{\Phi_4},
\qquad
\Phi_4 \;=\; q^2 + 1 \;=\; 10,
\]
and the chiral-sector pair
\[
u_{\mathrm{chiral}} \;=\; -2 \;\pm\; i\sqrt{\Phi_6},
\qquad
\Phi_6 \;=\; q^2 - q + 1 \;=\; 7,
\]
both with squared modulus equal to $k - 1 = 11 = p_{\mathrm{Ih}}$ (the
Ihara prime).  The squared imaginary parts are exactly the substrate
cyclotomic primitives $\Phi_4$ and $\Phi_6$, so the cyclotomic structure
of $W(3,3)$ is written directly into the directed photonic carrier.

Translating to phase angles:
\begin{equation}
\theta_{\mathrm{gauge}} = \arctan\!\sqrt{\Phi_4} \approx 72.45^\circ,
\qquad
\theta_{\mathrm{chiral}} = \pi - \arctan\!\bigl(\sqrt{\Phi_6}/\lambda\bigr)
                        \approx 127.09^\circ.
\end{equation}
These are the leading non-backtracking transport phases of the substrate.
In a single-photon dual-rail implementation, $\theta_{\mathrm{gauge}}$
and $\theta_{\mathrm{chiral}}$ govern the inter-edge interference fringes
visible after one Hashimoto step on the $480$-carrier, and supply a
\emph{falsifiable prediction} of the photonic runtime: any faithful
realisation of the substrate must exhibit these two phases as the
spectral content of one non-backtracking transport tick.


\subsection{Response and Operation Gate}

The post-measurement response layer is generated by the W33 two-graph odd triples.
If $M$ is the vertex-by-odd-triple incidence matrix, the imported bridge stack
verifies
\[
MM^{T}=320I+16J+4A,
\]
so adjacency, hence the feed-forward geometry, is recovered from parity incidence
rather than inserted by hand.  The same triangle complex has integral homology
$H_1\cong \mathbb{Z}^{81}$: its Smith certificate has rank $120$ and no torsion.

The current E8 operation gate is therefore exact but bounded.  The rank $81$ H1
module matches the E8 $\mathbb{Z}_3$ matter grades
\[
E_8=g_0(86)\oplus g_1(81)\oplus g_2(81),
\]
and the local operation pipeline verifies the grade rules on $8347$ nonzero bracket
terms.  It also pins $|g_1\times g_2|=81^2$ and $810$ nonzero $g_1\times g_1$
brackets, with $162$ firewall-filtered couplings.  This is the operation-level
gate into the E8 layer; it does not by itself prove a biological origin theorem.

\subsection{Protected Scheduler}

The runtime separation above becomes an executable compiler only after the
probabilistic, deterministic, protected, and classical handoffs are ordered.  The
current protected schedule has eight ticks:
\[
\begin{array}{ccl}
0&:&\text{projective carrier: }3^4=81\to 40,\\
1&:&\text{heralded fusion assembly: }p_{\mathrm{fusion}}=1/2,\;E/p=480,\\
2&:&\text{KLM primitive budget: }p_{\mathrm{KLM}}=1/4,\;E/p=960,\\
3&:&\text{CSS resource validation: }39+120+81=240,\\
4&:&\text{MBQC feed-forward: }4\text{ frame trits}=3^4=81,\\
5&:&\text{Steane}/\Phi_6\text{ protection: }[[240,81,3]]\to[[82320,81,\ge81]],\\
6&:&\text{classical selector commit: }2^{63}<3^{40}<2^{64},\\
7&:&\text{E8 } \mathbb{Z}_3\text{ operation gate: }8347\text{ checked brackets.}
\end{array}
\]
The tick count is not used as a numerological proof; it is the finite controller
rank needed to keep the handoff interfaces distinct.  Tick 3 is the bare W33
edge-qubit CSS code from the vertex and triangle checks.  Its distance is exactly
3, so tick 5 applies the Steane code~\cite{Steane1996} as a $\Phi_6=7$ inner block
three times.  The resulting code has parameter lower bound
[[82320,81,\(\ge 81\)]], hence it corrects at least \(40\) arbitrary faults, exactly
the number of W33 measurement trits committed at tick 6.

Thus the scheduler is the strict version of the architecture claim: probabilistic
photonic events are retried while still physical; once the graph resource is
accepted, MBQC randomness is deterministic Pauli-frame data; the protected code
absorbs the finite quantum fault budget; and only then is the 40-trit record exposed
as a classical selector.

\subsection{Fusion-Control Splice}

The scheduler's CSS validation tick can now be refined using the local
Csaszar theta packet.  The five input tori carry a local code
\[
5[[21,2,\ge 3]]=[[105,10,\ge 3]],
\]
so the local check rank is \(105-10=95\).  The five input modes contribute a
natural \(U(5)\) rank \(25\), and therefore
\[
95+25=120.
\]
This resolves the \(120\)-rank triangle term in tick 3.  The full edge-carrier
identity becomes
\[
95+25+39+81=240,
\]
where the last \(81\) is the protected H1 matter tail, not a stabilizer to be
removed.

The same refinement splits the physical carrier:
\[
105+135=240.
\]
Consequently the probabilistic runtime budgets split as
\[
p_{\mathrm{fusion}}=1/2:\qquad 210+270=480
\]
and
\[
p_{\mathrm{KLM}}=1/4:\qquad 420+540=960.
\]
Thus nondeterministic optical assembly is not being reinterpreted as a
deterministic gate.  It is a heralded/fusion budget whose accepted carrier is
then passed into the CSS and MBQC layers, in the same spirit as photonic fusion-based fault tolerance~\cite{Bartolucci2021}.  The deterministic layer is
still the H1-sized Pauli-frame lock \(3^4=81\), and the classical layer remains
the \(40\)-trit selector inside the \(64\)-bit envelope.

\subsection{Curved Product Handoff}

The protected finite scheduler is not, by itself, a four-dimensional continuum
geometry.  The current curved handoff is the almost-commutative product
\[
\Delta_{\mathrm{ext}}\otimes 1 + 1\otimes D_F^2,
\]
where the external factor is supplied by explicit simplicial four-dimensional
seeds.  The live certificates use \(CP^2_9\), with Betti profile
\((1,0,1,0,1)\), and \(K3_{16}\), with Betti profile \((1,0,22,0,1)\).  Thus the
protected H1 tail lifts through the external harmonic sectors as
\[
81\cdot 3=243,\qquad 81\cdot 24=1944.
\]
The barycentric refinement tower supplies the genuine four-dimensional scaling
family, with local density limits \(120/19\) and \(860/19\).  This is a product
handoff, not yet the final Einstein--Hilbert spectral-action asymptotic theorem.

\subsection{Curved Coefficient Extractor}

On the same curved tower, the first-moment sequence splits exactly into the
\(120\)-, \(6\)-, and \(1\)-mode channels.  Equivalently, its generating function
has the residue form
\[
\frac{A}{1-120z}+\frac{B}{1-6z}+\frac{C}{1-z}.
\]
The \(6\)-mode is the Einstein--Hilbert-like curvature channel.  For either
explicit seed, any three successive refinement samples recover
\[
c_6=12480,\qquad c_{\mathrm{EH}}=320,\qquad a_2=2240,
\]
with \(12480/39=320\).  The same roundtrip reconstructs
\[
D_F^2=\{0^{82},4^{320},10^{48},16^{30}\},
\]
the finite moments \((480,2240,17600)\), the internal
\(\mathrm{SRG}(40,12,2,4)\), and the promoted generator
\[
x=\sin^2\theta_W=\frac{3}{13}.
\]
So the photonic runtime now has a verified finite-to-curved coefficient
extractor.  The remaining open theorem is the smooth spectral-action limit, not
the finite coefficient roundtrip.

%======================================================================
\section{The \texorpdfstring{$W(3,3)$}{W(3,3)} Structural Encoding Theorem}
\label{sec:encoding}
%======================================================================

\begin{theorem}[Structural encoding of photonic universality in $W(3,3)$]
\label{thm:encoding}
Let $W(3,3)$ denote the symplectic polar space with parameters as in
\cref{prop:params}.  Then every cardinal number governing single-photon universal
quantum computation is a closed arithmetic expression in the $W(3,3)$ parameters.
\end{theorem}

\begin{center}
\renewcommand{\arraystretch}{1.18}
\begin{tabular}{@{}p{0.44\linewidth}cc@{}}
\toprule
\textbf{Physical / Computational Quantity} & \textbf{Value} & \textbf{$W(3,3)$ formula} \\
\midrule
Qubit dimension (photon polarisation states)       & $2$     & $\lambda = 2$ \\
Physical photon helicity states                    & $2$     & $\lambda = 2$ \\
Spinor (Weyl) dimension of the photon              & $2$     & $\lambda = 2$ \\
Dual-rail modes per logical qubit                  & $2$     & $\lambda = 2$ \\
Smallest UTM state count (Wolfram--Smith)          & $3$     & $q = 3$ \\
Pauli generators of $\mathrm{SU}(2)$              & $3$     & $q = 3$ \\
Wave plates for universal $\mathrm{SU}(2)$         & $3$     & $q = 3$ \\
Qutrit Hilbert space dimension ($\CC^3$)           & $3$     & $q = 3$ \\
OAM modes $\ell\in\{-1,0,1\}$ for photonic qutrit & $3$     & $q = 3$ \\
Spacetime dimensions                               & $4$     & $\mu = 4$ \\
KLM ancilla modes (elementary gadget)              & $4$     & $\mu = 4$ \\
$\mathrm{SU}(2)$ rank                             & $1$     & $q - \lambda = 1$ \\
SM gauge generators ($8+3+1$)                      & $12$    & $k = 12$ \\
Two-qutrit Clifford generators                     & $12$    & $k = 12$ \\
Vertices / photons in $W(3,3)$ cluster             & $40$    & $v = 40$ \\
Edges / entangling operations in cluster           & $240$   & $E = 240$ \\
Expected Type-II fusion attempts                   & $480$   & $E/(\lambda/\mu)=2E$ \\
Primitive KLM attempts for all edges               & $960$   & $E/(1/\mu)=\mathrm{tr}(A^3)$ \\
Two-qutrit Pauli-frame states                      & $81$    & $q^4$ \\
Single-run measurement word                        & $40$ trits & $2^{63}<3^{40}<2^{64}$ \\
Toric-code logical qubits on a torus               & $2$     & $\lambda = 2$ \\
Toric-code ground-state degeneracy on a torus      & $4$     & $\mu = 4$ \\
H1 operation module rank                           & $81$    & $q^4$ \\
Verified E8 Z3 bracket terms                       & $8347$  & pipeline certificate \\
$g_1\times g_2$ operation pairs                    & $6561$  & $q^8$ \\
$g_1\times g_1$ nonzero brackets                  & $810$   & $10q^4$ \\
Protected CSS code                                 & $[[82320,81,\ge81]]$ & $[240\Phi_6^3, q^4, \ge q^4]$ \\
Runtime scheduler ticks                            & $8$     & $2^q$ \\
Clifford automorphism group order                  & $51840$ & $|\Sp(4,\FF_3)|$ \\
Positive eigenvalue / qubit eigenvalue             & $2$     & $r = 2$ \\
Laplacian gap / cluster propagation bandwidth      & $10$    & $k - r = 10$ \\
Multiplicity of $r$ (Leech lattice kissing dim.)   & $24$    & $f = 24$ \\
\bottomrule
\end{tabular}
\end{center}

\begin{proof}[Proof of \cref{thm:encoding}]
The entries $\lambda=2$, $\mu=4$, $q=3$, $k=12$, $v=40$, $E=240$ follow from
\cref{prop:params}.  The identification $\lambda=2$ as qubit dimension follows from
$\lambda$ counting common neighbours, equal to the number of independent photon
polarisation modes.  The identification $q=3$ as qutrit dimension follows from
\cref{thm:dioph} and \cref{def:pauli3}.  The identification $k=12$ as Clifford
generator count follows from \cref{thm:phasespace}.  The Laplacian gap $k-r=10$ is
the graph-theoretic analogue of the mixing bandwidth.  The automorphism group order
$51840 = |\Sp(4,\FF_3)|$ is verified computationally.
\end{proof}

\begin{corollary}[Compression of physical law]
The parameters $(q=3, k=12, \lambda=2, \mu=4, v=40, E=240)$ can be encoded in fewer
than $64$ bits (Kolmogorov complexity $K(W(3,3)) < 64$), whereas independent
specification of the full Standard Model plus quantum optics parameter set requires
$\gtrsim 260$ bits.  The $W(3,3)$ skeleton achieves a $6.5\times$ compression of the
fundamental constants governing photonic universality.
\end{corollary}

%======================================================================
\section{Conclusions}
\label{sec:conclusions}
%======================================================================

We have demonstrated, from three complementary perspectives---circuit-based linear
optics, the KLM non-deterministic gate protocol, and measurement-based cluster
computation---that a single photon is the fundamental unit of universal quantum
computation.  The central new contribution of this paper is to ground this claim in
the symplectic polar geometry of $W(3,3)$: not merely a graph with four parameters,
but a classical incidence geometry over the ternary field $\GF{3}$, whose existence
is forced by the arithmetic identity $3! = 2 \cdot 3$.

Key findings:
\begin{itemize}[noitemsep]
  \item \textbf{The field $\GF{3}$ is uniquely selected} by the Diophantine equation
    $q! = 2q$; no other field yields this coincidence.
  \item \textbf{$W(3,3)$ is a symplectic polar space} on $\PG{3}{\GF{3}}$, whose 40
    points are the non-zero two-qutrit Pauli observables and whose automorphism group
    is $\Sp(4,\FF_3)$---the two-qutrit Clifford group.
  \item \textbf{Single-qubit universality} via passive optics; qubit dimension
    $\lambda=2$ and wave-plate count $q=3$ are $W(3,3)$ parameters.
  \item \textbf{Photonic qutrits} from three-mode photons; Clifford$+T_3$ universal,
    with symmetry group $\Sp(4,\FF_3)$ of order $51840$.
  \item \textbf{KLM universality} via ancilla count $\mu=4$.
  \item \textbf{MBQC universality} via $W(3,3)$ cluster state: 40 photons, 240
    entangling operations, Laplacian bandwidth $k - r = 10$.
  \item \textbf{Runtime separation:} probabilistic optical assembly feeds a
    deterministic MBQC Pauli frame and a 40-trit classical measurement record.
  \item \textbf{Protected scheduler:} the executable handoff has eight ticks, from
    projective carrier through heralded fusion, CSS validation, MBQC feed-forward,
    Steane/$\Phi_6$ protection, 64-bit-class selector commit, and E8 operation.
  \item \textbf{Fusion-control splice:} the \(105+135=240\) theta/transport carrier
    refines the runtime budgets as \(210+270=480\) fusion attempts and
    \(420+540=960\) KLM primitives, while tick 3 refines to
    \(95+25+39+81=240\).
  \item \textbf{Structural encoding theorem:} the listed cardinal numbers of
    photonic universality are closed arithmetic expressions in the $W(3,3)$
    parameter table.
\end{itemize}

Future directions include: loss-tolerant encodings exploiting the $E_6$ Weyl group
symmetry $\PSp(4,3) \cong W(E_6)^+$; boson-sampling complexity connected to the
permanent of the $W(3,3)$ adjacency matrix; fault-tolerant photonic codes using
$|\Aut(W(3,3))| = 51840$ as a symmetry resource; and higher-dimensional analogues
$W(2n-1, p)$ for prime $p > 3$.

%======================================================================
\begin{thebibliography}{99}

\bibitem{KLM2001}
  E.~Knill, R.~Laflamme, and G.~J.~Milburn,
  ``A scheme for efficient quantum computation with linear optics,''
  \textit{Nature} \textbf{409}, 46--52 (2001).

\bibitem{Reck1994}
  M.~Reck, A.~Zeilinger, H.~J.~Bernstein, and P.~Bertani,
  ``Experimental realization of any discrete unitary operator,''
  \textit{Phys.~Rev.~Lett.} \textbf{73}, 58 (1994).

\bibitem{Raussendorf2001}
  R.~Raussendorf and H.~J.~Briegel,
  ``A one-way quantum computer,''
  \textit{Phys.~Rev.~Lett.} \textbf{86}, 5188 (2001).

\bibitem{Browne2005}
  D.~E.~Browne and T.~Rudolph,
  ``Resource-efficient linear optical quantum computation,''
  \textit{Phys.~Rev.~Lett.} \textbf{95}, 010501 (2005).

\bibitem{Bartolucci2021}
  S.~Bartolucci \textit{et al.},
  ``Fusion-based quantum computation,''
  arXiv:2101.09310 (2021).

\bibitem{Steane1996}
  A.~M.~Steane,
  ``Multiple-particle interference and quantum error correction,''
  \textit{Proc.~R.~Soc.~Lond.~A} \textbf{452}, 2551--2577 (1996).

\bibitem{CalderbankShor1996}
  A.~R.~Calderbank and P.~W.~Shor,
  ``Good quantum error-correcting codes exist,''
  \textit{Phys.~Rev.~A} \textbf{54}, 1098--1105 (1996).

\bibitem{Kok2007}
  P.~Kok \textit{et al.},
  ``Linear optical quantum computing with photonic qubits,''
  \textit{Rev.~Mod.~Phys.} \textbf{79}, 135 (2007).

\bibitem{Smith2007}
  A.~Smith,
  ``Universality of Wolfram's 2, 3 Turing Machine,''
  \textit{Complex Systems} \textbf{19} (2007).

\bibitem{BK2005}
  S.~Bravyi and A.~Kitaev,
  ``Universal quantum computation with ideal Clifford gates and noisy ancillas,''
  \textit{Phys.~Rev.~A} \textbf{71}, 022316 (2005).

\bibitem{Erhard2018}
  M.~Erhard \textit{et al.},
  ``Twisted photons: new quantum perspectives in high dimensions,''
  \textit{Light Sci.~Appl.} \textbf{7}, 17146 (2018).

\bibitem{Tits1974}
  J.~Tits,
  \textit{Buildings of Spherical Type and Finite BN-Pairs},
  Lecture Notes in Mathematics 386, Springer, 1974.

\bibitem{Cameron1991}
  P.~J.~Cameron,
  \textit{Projective and Polar Spaces},
  QMW Maths Notes 13, 1991.

\bibitem{Brendel1999}
  J.~Brendel, N.~Gisin, W.~Tittel, and H.~Zbinden,
  \textit{Pulsed Energy-Time Entangled Twin-Photon Source for Quantum
  Communication},
  Phys.\ Rev.\ Lett.\ \textbf{82}, 2594 (1999).

\bibitem{Marcikic2002}
  I.~Marcikic, H.~de Riedmatten, W.~Tittel, V.~Scarani, H.~Zbinden, and
  N.~Gisin,
  \textit{Time-bin entangled qubits for quantum communication created by
  femtosecond pulses},
  Phys.\ Rev.\ A \textbf{66}, 062308 (2002).

\bibitem{Bennett1993}
  C.~H.~Bennett, G.~Brassard, C.~Cr\'epeau, R.~Jozsa, A.~Peres, and
  W.~K.~Wootters,
  \textit{Teleporting an Unknown Quantum State via Dual Classical and
  Einstein--Podolsky--Rosen Channels},
  Phys.\ Rev.\ Lett.\ \textbf{70}, 1895 (1993).

\bibitem{Vlasov2025}
  A.~Yu.~Vlasov,
  \textit{Witting polytope, quantum communication, and the symplectic polar
  space $W(3,3)$},
  preprint, 2025; see also W(3,3) For Everyone, Parts MCLXIII--MCCIII.

\end{thebibliography}

\end{document}
